% Options for packages loaded elsewhere
\PassOptionsToPackage{unicode}{hyperref}
\PassOptionsToPackage{hyphens}{url}
\PassOptionsToPackage{dvipsnames,svgnames,x11names}{xcolor}
\documentclass[
  11pt,
]{article}
\usepackage{xcolor}
\usepackage[margin=25mm]{geometry}
\usepackage{amsmath,amssymb}
\setcounter{secnumdepth}{-\maxdimen} % remove section numbering
\usepackage{iftex}
\ifPDFTeX
  \usepackage[T1]{fontenc}
  \usepackage[utf8]{inputenc}
  \usepackage{textcomp} % provide euro and other symbols
\else % if luatex or xetex
  \usepackage{unicode-math} % this also loads fontspec
  \defaultfontfeatures{Scale=MatchLowercase}
  \defaultfontfeatures[\rmfamily]{Ligatures=TeX,Scale=1}
\fi
\usepackage{lmodern}
\ifPDFTeX\else
  % xetex/luatex font selection
\fi
% Use upquote if available, for straight quotes in verbatim environments
\IfFileExists{upquote.sty}{\usepackage{upquote}}{}
\IfFileExists{microtype.sty}{% use microtype if available
  \usepackage[]{microtype}
  \UseMicrotypeSet[protrusion]{basicmath} % disable protrusion for tt fonts
}{}
\makeatletter
\@ifundefined{KOMAClassName}{% if non-KOMA class
  \IfFileExists{parskip.sty}{%
    \usepackage{parskip}
  }{% else
    \setlength{\parindent}{0pt}
    \setlength{\parskip}{6pt plus 2pt minus 1pt}}
}{% if KOMA class
  \KOMAoptions{parskip=half}}
\makeatother
\usepackage{color}
\usepackage{fancyvrb}
\newcommand{\VerbBar}{|}
\newcommand{\VERB}{\Verb[commandchars=\\\{\}]}
\DefineVerbatimEnvironment{Highlighting}{Verbatim}{commandchars=\\\{\}}
% Add ',fontsize=\small' for more characters per line
\newenvironment{Shaded}{}{}
\newcommand{\AlertTok}[1]{\textcolor[rgb]{1.00,0.00,0.00}{\textbf{#1}}}
\newcommand{\AnnotationTok}[1]{\textcolor[rgb]{0.38,0.63,0.69}{\textbf{\textit{#1}}}}
\newcommand{\AttributeTok}[1]{\textcolor[rgb]{0.49,0.56,0.16}{#1}}
\newcommand{\BaseNTok}[1]{\textcolor[rgb]{0.25,0.63,0.44}{#1}}
\newcommand{\BuiltInTok}[1]{\textcolor[rgb]{0.00,0.50,0.00}{#1}}
\newcommand{\CharTok}[1]{\textcolor[rgb]{0.25,0.44,0.63}{#1}}
\newcommand{\CommentTok}[1]{\textcolor[rgb]{0.38,0.63,0.69}{\textit{#1}}}
\newcommand{\CommentVarTok}[1]{\textcolor[rgb]{0.38,0.63,0.69}{\textbf{\textit{#1}}}}
\newcommand{\ConstantTok}[1]{\textcolor[rgb]{0.53,0.00,0.00}{#1}}
\newcommand{\ControlFlowTok}[1]{\textcolor[rgb]{0.00,0.44,0.13}{\textbf{#1}}}
\newcommand{\DataTypeTok}[1]{\textcolor[rgb]{0.56,0.13,0.00}{#1}}
\newcommand{\DecValTok}[1]{\textcolor[rgb]{0.25,0.63,0.44}{#1}}
\newcommand{\DocumentationTok}[1]{\textcolor[rgb]{0.73,0.13,0.13}{\textit{#1}}}
\newcommand{\ErrorTok}[1]{\textcolor[rgb]{1.00,0.00,0.00}{\textbf{#1}}}
\newcommand{\ExtensionTok}[1]{#1}
\newcommand{\FloatTok}[1]{\textcolor[rgb]{0.25,0.63,0.44}{#1}}
\newcommand{\FunctionTok}[1]{\textcolor[rgb]{0.02,0.16,0.49}{#1}}
\newcommand{\ImportTok}[1]{\textcolor[rgb]{0.00,0.50,0.00}{\textbf{#1}}}
\newcommand{\InformationTok}[1]{\textcolor[rgb]{0.38,0.63,0.69}{\textbf{\textit{#1}}}}
\newcommand{\KeywordTok}[1]{\textcolor[rgb]{0.00,0.44,0.13}{\textbf{#1}}}
\newcommand{\NormalTok}[1]{#1}
\newcommand{\OperatorTok}[1]{\textcolor[rgb]{0.40,0.40,0.40}{#1}}
\newcommand{\OtherTok}[1]{\textcolor[rgb]{0.00,0.44,0.13}{#1}}
\newcommand{\PreprocessorTok}[1]{\textcolor[rgb]{0.74,0.48,0.00}{#1}}
\newcommand{\RegionMarkerTok}[1]{#1}
\newcommand{\SpecialCharTok}[1]{\textcolor[rgb]{0.25,0.44,0.63}{#1}}
\newcommand{\SpecialStringTok}[1]{\textcolor[rgb]{0.73,0.40,0.53}{#1}}
\newcommand{\StringTok}[1]{\textcolor[rgb]{0.25,0.44,0.63}{#1}}
\newcommand{\VariableTok}[1]{\textcolor[rgb]{0.10,0.09,0.49}{#1}}
\newcommand{\VerbatimStringTok}[1]{\textcolor[rgb]{0.25,0.44,0.63}{#1}}
\newcommand{\WarningTok}[1]{\textcolor[rgb]{0.38,0.63,0.69}{\textbf{\textit{#1}}}}
\usepackage{longtable,booktabs,array}
\usepackage{calc} % for calculating minipage widths
% Correct order of tables after \paragraph or \subparagraph
\usepackage{etoolbox}
\makeatletter
\patchcmd\longtable{\par}{\if@noskipsec\mbox{}\fi\par}{}{}
\makeatother
% Allow footnotes in longtable head/foot
\IfFileExists{footnotehyper.sty}{\usepackage{footnotehyper}}{\usepackage{footnote}}
\makesavenoteenv{longtable}
\setlength{\emergencystretch}{3em} % prevent overfull lines
\providecommand{\tightlist}{%
  \setlength{\itemsep}{0pt}\setlength{\parskip}{0pt}}
\emergencystretch=3em
% Disable common ligatures so copy/search text retains the literal fi/fl
% sequences across PDF readers and text extractors.
\defaultfontfeatures{Ligatures=NoCommon}
\IfFontExistsTF{DejaVu Serif}
  {\setmainfont{DejaVu Serif}[Ligatures=NoCommon]}
  {\setmainfont{Cambria}[Ligatures=NoCommon]}
\IfFontExistsTF{DejaVu Sans Mono}
  {\setmonofont{DejaVu Sans Mono}[Scale=MatchLowercase,Ligatures=NoCommon]}
  {\setmonofont{Cascadia Mono}[Scale=MatchLowercase,Ligatures=NoCommon]}
\IfFontExistsTF{DejaVu Sans}
  {\newfontfamily\symbolfont{DejaVu Sans}}
  {\newfontfamily\symbolfont{Segoe UI Symbol}}
\usepackage{xurl}
\usepackage{adjustbox}
\usepackage{fvextra}
% Long highlighted directory trees must wrap rather than silently leave the page.
\fvset{breaklines=true,breakanywhere=true,breaknonspaceingroup=true}
\usepackage{newunicodechar}
\newunicodechar{ℓ}{\ensuremath{\ell}}
\newunicodechar{ℂ}{\ensuremath{\mathbb{C}}}
\newunicodechar{⊬}{\ensuremath{\nvdash}}
\newunicodechar{★}{\mbox{\symbolfont\char"2605}}
\newunicodechar{𝔻}{\ensuremath{\mathbb{D}}}
\newunicodechar{ℍ}{\ensuremath{\mathbb{H}}}
\newunicodechar{⪯}{\ensuremath{\preceq}}
\newunicodechar{⪰}{\ensuremath{\succeq}}
\newunicodechar{⁅}{\mbox{\symbolfont\char"2045}}
\newunicodechar{⁆}{\mbox{\symbolfont\char"2046}}
\newunicodechar{‖}{\ensuremath{\Vert}}
\newunicodechar{⇒}{\ensuremath{\Rightarrow}}
\newunicodechar{⇏}{\ensuremath{\nRightarrow}}
\newunicodechar{∇}{\ensuremath{\nabla}}
\newunicodechar{∉}{\ensuremath{\notin}}
\newunicodechar{∘}{\ensuremath{\circ}}
\newunicodechar{∧}{\ensuremath{\land}}
\newunicodechar{⊂}{\ensuremath{\subset}}
\newunicodechar{⊆}{\ensuremath{\subseteq}}
\newunicodechar{⊗}{\ensuremath{\otimes}}
\newunicodechar{♯}{\ensuremath{\sharp}}
\newunicodechar{ᵇ}{\ensuremath{^{\mathrm b}}}
\newunicodechar{ᵏ}{\ensuremath{^{\mathrm k}}}
\newunicodechar{ₙ}{\ensuremath{_{\mathrm n}}}
\newunicodechar{ᵀ}{\ensuremath{^{\mathsf T}}}
\AtBeginDocument{\Urlmuskip=0mu plus 2mu\urlstyle{tt}}
\setlength{\tabcolsep}{4pt}
\usepackage{bookmark}
\IfFileExists{xurl.sty}{\usepackage{xurl}}{} % add URL line breaks if available
\urlstyle{same}
\hypersetup{
  pdftitle={Finite Quantum Mechanics},
  pdfauthor={Jeremy H. Carroll},
  colorlinks=true,
  linkcolor={Maroon},
  filecolor={Maroon},
  citecolor={Blue},
  urlcolor={Blue},
  pdfcreator={LaTeX via pandoc}}

\title{Finite Quantum Mechanics}
\author{Jeremy H. Carroll}
\date{2026-10-14}

\begin{document}
\makeatletter
\begin{center}
{\LARGE\bfseries \@title\par}
\vspace{0.45em}
{\large Can Quantum Mechanics Be Presented on a Finite Carrier --- and If Not,
Exactly Where Does It Break?\par}
\vspace{0.8em}
{\normalsize \@author\par}
\vspace{0.5em}
{\small \@date\par}
\vspace{0.2em}
{\small Version 0.55.0\par}
\vspace{0.3em}
{\footnotesize Creative Commons Attribution 4.0 International (CC-BY-4.0)\par}
\end{center}
\makeatother

{
\hypersetup{linkcolor=}
\setcounter{tocdepth}{2}
\tableofcontents
}
\subsection{Abstract}\label{abstract}

We study finite-dimensional quantum models over specified rational and
algebraic scalar fields, distinguishing a finite presentation from a
finite set of states. The principal results are exact constructions and
obstructions under explicit assumptions on states, effects, composition,
dynamics, and readout. A normalized clock seed gives a constructive
finite Weyl intertwiner. Equal population reports can fail to preserve
answers to a common future experiment. An unconditional
measure-and-prepare transfer of an unknown qubit has average survival
probability at most three quarters on four specified preparations, with
an attaining strategy. Finite Bell tables distinguish support
compatibility from the existence of a positive common joint
distribution, while central-sector weights constrain purification in a
specified direct-sum tensor model.

Exact finite examples, written arguments, and selected formalizations
accompany one another with separate scopes. A positive global joint law
implies the CHSH bound \(|S|\leq 2\); a supplied rational quantum model
reaches \(S=14/5\); and a supplied tensor-state-observable model over
\(\mathbb Q(\sqrt 2)\) attains \(S=2\sqrt 2\) and has an exact
sum-of-squares upper certificate. These results do not derive the
quantum state space, Born pairing, tensor composition, or continuous
dynamics from primitive difference. They identify which supplied
structures each conclusion uses and which reconstruction questions
remain open.

\subsection{Guide to the paper}\label{guide-to-the-paper}

The main text is organized by mathematical questions. Source provenance
and reproduction details accompany the results without determining their
truth.

\begin{longtable}[]{@{}
  >{\raggedright\arraybackslash}p{(\linewidth - 4\tabcolsep) * \real{0.3333}}
  >{\raggedright\arraybackslash}p{(\linewidth - 4\tabcolsep) * \real{0.3333}}
  >{\raggedright\arraybackslash}p{(\linewidth - 4\tabcolsep) * \real{0.3333}}@{}}
\toprule\noalign{}
\begin{minipage}[b]{\linewidth}\raggedright
Question
\end{minipage} & \begin{minipage}[b]{\linewidth}\raggedright
Location
\end{minipage} & \begin{minipage}[b]{\linewidth}\raggedright
Evidential boundary
\end{minipage} \\
\midrule\noalign{}
\endhead
\bottomrule\noalign{}
\endlastfoot
Which finite and algebraic structures are being assumed? & §1 & Declared
carrier and scalar data are not derived conclusions \\
What counts as a proof or finite certificate? & §§2--3 & Evidence is
matched to the quantifiers of the statement \\
Which mathematical results are established? & §§4 and 11 & Each theorem
keeps its scalar, dimension, positivity, and composition premises \\
Where does the proposed reconstruction remain open? & §§5, 7, and 9 & A
missing selection argument is not an impossibility theorem \\
How can the computations be reproduced? & §§6, 10, and 12 & Validation
and serialization checks do not establish physical predicates that
fields only name \\
What do the carrier and measurement examples establish? & §13 &
Conditional operational separations; no individual-outcome selection
law \\
\end{longtable}

\subsection{1. Finite presentations, supplied structure, and target
examples}\label{finite-presentations-supplied-structure-and-target-examples}

A finite-dimensional quantum model need not have finitely many states.
This paper distinguishes four notions:

\begin{enumerate}
\def\labelenumi{\arabic{enumi}.}
\tightlist
\item
  a finite graph or cell complex;
\item
  a finite presentation by exact algebraic data;
\item
  a finite-dimensional vector space, which normally has infinitely many
  normalized states; and
\item
  a finite list of preparations, effects, or experiments.
\end{enumerate}

On a declared directed graph the primitive coboundary is

\[\protect\adjustbox{max width=\linewidth}{\(\displaystyle (d_0f)(u,v)=f(v)-f(u).\)}\]

For an edge cochain \(x\), a weighted quotient such as

\[\protect\adjustbox{max width=\linewidth}{\(\displaystyle \Phi_1([x])=\inf_{\alpha}\|x-d_0\alpha\|_{1,w}\)}\]

measures the residual after the admitted vertex repairs. This
construction is mathematically useful, but it does not by itself select
complex scalars, positivity, inner products, tensor composition, the
Born pairing, or continuous time. Those choices must be supplied or
derived separately.

\subsubsection{1.1 An eight-rung
inventory}\label{an-eight-rung-inventory}

The eight rungs below are an assumption audit, not an eight-step
derivation of quantum mechanics.

\begin{longtable}[]{@{}
  >{\raggedright\arraybackslash}p{(\linewidth - 6\tabcolsep) * \real{0.2500}}
  >{\raggedright\arraybackslash}p{(\linewidth - 6\tabcolsep) * \real{0.2500}}
  >{\raggedright\arraybackslash}p{(\linewidth - 6\tabcolsep) * \real{0.2500}}
  >{\raggedright\arraybackslash}p{(\linewidth - 6\tabcolsep) * \real{0.2500}}@{}}
\toprule\noalign{}
\begin{minipage}[b]{\linewidth}\raggedright
Rung
\end{minipage} & \begin{minipage}[b]{\linewidth}\raggedright
Established once the inputs are supplied
\end{minipage} & \begin{minipage}[b]{\linewidth}\raggedright
Additional supplied structure
\end{minipage} & \begin{minipage}[b]{\linewidth}\raggedright
Remaining boundary
\end{minipage} \\
\midrule\noalign{}
\endhead
\bottomrule\noalign{}
\endlastfoot
1. Difference & Finite differences, equality tests, and graph
coboundaries & A carrier and equality relation & No preferred norm,
topology, or physical geometry is selected \\
2. Quotient repair & Exact finite quotient values and
primal\allowbreak{}/\allowbreak{}dual certificates & Coefficient space,
weights, and repair image & The quotient does not choose a quantum state
space \\
3. State--effect pairing & Positive normalized weights and conditional
identities in stated models & Positive states, effects, and a pairing &
The Born rule is not derived from difference \\
4. Phase & Relative phases and phase-sensitive interference on a
supplied complex or cyclotomic carrier & Complex conjugation and a phase
group & Population data alone do not select the phase fiber \\
5. Composition & Product states, entangled states, partial traces, and
Bell tables in a supplied tensor model & Tensor composition and
positivity & Tensor products are not forced by graph products \\
6. Reversible dynamics & Exact finite unitaries and rational Cayley
steps & An inner product, reversible maps, and a time step & Continuous
time and a generator require topology and interpolation \\
7. Measurement and channels & POVM probabilities, conditional updates,
exact channels, and finite dilations & Effect algebra, instruments, and
completely positive maps & Unitary dilation does not select one realized
outcome \\
8. Representation boundary & Finite frame constraints, contextuality
witnesses, and finite-dimensional commutator obstructions & The
hypotheses of Gleason\allowbreak{}/\allowbreak{}Busch-type
representation theorems & Restricted algebraic ray sets do not
automatically extend to complete projective geometry \\
\end{longtable}

Two elementary cautions are important. First, the difference and
equality relation do not select the norm, topology, or physical geometry
needed later; they do not exclude every metric construction. Second,
projective state space is

\[\protect\adjustbox{max width=\linewidth}{\(\displaystyle \mathbb P(\mathcal H)
=(\mathcal H\setminus\{0\})/\mathbb C^\times
\cong S(\mathcal H)/U(1),\)}\]

where \(S(\mathcal H)\) is the unit sphere. The group quotient
\(\mathbb C^\times/U(1)\) instead records positive scale and is not the
projectivization of the state space.

\subsubsection{1.2 Twenty-two finite
targets}\label{twenty-two-finite-targets}

The target list tests exact finite presentations against familiar
structures. Each row states the object actually considered; it is not a
claim that the corresponding general physical theory has been
reconstructed.

\begin{longtable}[]{@{}
  >{\raggedright\arraybackslash}p{(\linewidth - 2\tabcolsep) * \real{0.5000}}
  >{\raggedright\arraybackslash}p{(\linewidth - 2\tabcolsep) * \real{0.5000}}@{}}
\toprule\noalign{}
\begin{minipage}[b]{\linewidth}\raggedright
Target
\end{minipage} & \begin{minipage}[b]{\linewidth}\raggedright
Finite object and present scope
\end{minipage} \\
\midrule\noalign{}
\endhead
\bottomrule\noalign{}
\endlastfoot
1 & Primitive difference and zero self-difference on a declared
carrier \\
2 & A diamond-shaped 2-complex with exact cyclotomic holonomy data \\
3 & A three-path amplitude fixture with vanishing Sorkin third-order
term {[}7{]} under a supplied quadratic readout \\
4 & Rational Cayley transforms in \(\mathfrak{so}(8,\mathbb Q)\) and
exact orthogonality checks \\
5 & Rational \(S_3\) braid relations and recoupling-shaped examples; no
general pentagon theorem \\
6 & Cabello's 18-ray parity system, its 146-vertex rational frame
polytope, and exact positivity filtering \\
7 & Signed permutation unitaries on a finite layered model \\
8 & A rational CHSH model with \(S=14/5\) and an exact
positive-extension obstruction \\
9 & The three-qubit bit-flip code, whose distance is one against
unrestricted Pauli errors \\
10 & Four-branch teleportation identities over exact Gaussian-rational
data \\
11 & Superdense-coding identities on a supplied Bell pair and gate
set \\
12 & A four-qubit \([[4,2,2]]\) stabilizer example with exact syndrome
and distance checks \\
13 & Prime-\(5\) discrete Wigner phase-point operators and selected
covariance\allowbreak{}/\allowbreak{}negativity witnesses \\
14 & The Shor \([[9,1,3]]\) stabilizer construction and single-qubit
error syndromes \\
15 & Four-point phase-estimation identities over a specified cyclotomic
carrier \\
16 & A finite toric\allowbreak{}/\allowbreak{}surface-code incidence
model and its declared anyon relations \\
17 & Four-item Grover search as two exact Householder reflections \\
18 & Finite Clifford conjugation identities and a scoped
transversal-gate limitation {[}25{]} \\
19 & The finite commutator trace obstruction and a constructive finite
Weyl representation \\
20 & Exact Richardson weights and the variance cost of finite zero-noise
extrapolation \\
21 & A finite quantum walk on a supplied Cayley graph with exact norm
preservation \\
22 & Exact mutually unbiased-basis tomography in dimensions three and
four, with positivity checked separately from linear reconstruction \\
\end{longtable}

The detailed target definitions, implementation addresses, and
formalization scope are in the source-notes supplement. Three results
deserve emphasis here: the Bell separation in §11.4, the Cabello census
in §11.5, and the following finite Weyl construction.

\subsubsection{1.3 Constructive finite Weyl equivalence from an anchored
seed}\label{constructive-finite-weyl-equivalence-from-an-anchored-seed}

Let \(N \ge 1\) and let \(K \subseteq \mathbb C\) be a
conjugation-stable subfield containing a primitive \(N\)th root of unity
\(\omega\). Equip \(K^N\) with the standard Hermitian pairing. Supply
unitary \(N \times N\) matrices \(X', Z'\) over \(K\) satisfying the
finite Weyl relations {[}4{]}

\[\protect\adjustbox{max width=\linewidth}{\(\displaystyle Z'X'=\omega X'Z',\qquad (X')^N=I,\)}\]

and a normalized vector \(v \in K^N\) satisfying

\[\protect\adjustbox{max width=\linewidth}{\(\displaystyle Z'v=v,\qquad v^\dagger v=1.\)}\]

Define the canonical shift and clock by \(X e_j = e_{j+1 \pmod N}\) and
\(Z e_j = \omega^j e_j\), for \(0 \le j < N\). Then

\[\protect\adjustbox{max width=\linewidth}{\(\displaystyle W = \begin{pmatrix} v & X'v & \cdots & (X')^{N-1}v \end{pmatrix}\)}\]

is unitary and satisfies \(X'W=WX\) and \(Z'W=WZ\). The joint commutant
of \(X'\) and \(Z'\) consists of scalar matrices. Any two unitary
intertwiners between the fixed pairs differ by multiplication by a
unit-modulus scalar in \(K\).

\emph{Proof.} Induction in the Weyl relation gives
\(Z'(X')^j v = \omega^j (X')^j v\). These \(N\) eigenvalues are
distinct. Since \(Z'\) is unitary, their eigenvectors are orthogonal,
and unitarity of \(X'\) makes all \(N\) orbit vectors have norm one. The
explicitly \(N\)-dimensional carrier makes them a basis. The period of
\(X'\) closes the shift relation, and the displayed eigenvalue equation
gives the clock relation. A matrix commuting with the canonical clock is
diagonal; commuting with the shift makes all its diagonal entries equal.
Conjugating by \(W\) establishes the commutant statement, and applying
it to the ratio of two intertwiners establishes uniqueness up to scalar.

The normalized eigenvalue-one seed and scalar field are supplied inputs.
This result does not find or normalize such a seed over every number
field. If a supplied seed has eigenvalue \(\omega^r\) instead, replace
it by \((X')^{-r}v\) before using this construction. On a larger carrier
the orbit spans a cyclic invariant subspace, not necessarily the whole
carrier; no scalar-commutant conclusion for the larger carrier follows.

\subsection{2. Standards of evidence}\label{standards-of-evidence}

The mathematical statement, its premises, its implementation and its
reproduction record are separate. We use the following evidence classes
throughout the paper:

\begin{longtable}[]{@{}
  >{\raggedright\arraybackslash}p{(\linewidth - 4\tabcolsep) * \real{0.3333}}
  >{\raggedright\arraybackslash}p{(\linewidth - 4\tabcolsep) * \real{0.3333}}
  >{\raggedright\arraybackslash}p{(\linewidth - 4\tabcolsep) * \real{0.3333}}@{}}
\toprule\noalign{}
\begin{minipage}[b]{\linewidth}\raggedright
Class
\end{minipage} & \begin{minipage}[b]{\linewidth}\raggedright
What it supports
\end{minipage} & \begin{minipage}[b]{\linewidth}\raggedright
What it does not support
\end{minipage} \\
\midrule\noalign{}
\endhead
\bottomrule\noalign{}
\endlastfoot
Written proof & The quantified statement under the assumptions used in
its argument & A claim of Lean verification or physical realization \\
Scoped Lean proof & The actual elaborated declaration, including its
input hypotheses and coefficient types & A broader theorem suggested by
a filename, comment or surrounding prose \\
Exact finite computation & An identity, enumeration or certificate
replay on its declared input domain & An unrestricted theorem from a
sample of instances \\
Numerical demonstration & A reported approximation with its algorithm
and tolerance & Exact zero, exact positivity or exact optimality without
additional certificates \\
Imported theorem or conditional model & A standard theorem with its
hypotheses, or a result conditional on supplied structures & Derivation
of those hypotheses or structures from finite history data \\
Historical observation / unreproduced claim & What an identified earlier
source or execution record reported & A fresh successful run or
endorsement of every present source claim \\
\end{longtable}

A finite computation can decide a finite theorem completely, and a
checked linear basis can establish an identity for its entire
finite-dimensional space. Conversely, a general written proof remains a
proof even when no Lean mechanization is supplied. Neither direction
should be collapsed into the ledger's software status labels.

The most important conclusions have these boundaries:

\begin{longtable}[]{@{}
  >{\raggedright\arraybackslash}p{(\linewidth - 4\tabcolsep) * \real{0.3333}}
  >{\raggedright\arraybackslash}p{(\linewidth - 4\tabcolsep) * \real{0.3333}}
  >{\raggedright\arraybackslash}p{(\linewidth - 4\tabcolsep) * \real{0.3333}}@{}}
\toprule\noalign{}
\begin{minipage}[b]{\linewidth}\raggedright
Conclusion
\end{minipage} & \begin{minipage}[b]{\linewidth}\raggedright
Premises and supporting argument
\end{minipage} & \begin{minipage}[b]{\linewidth}\raggedright
Current scope
\end{minipage} \\
\midrule\noalign{}
\endhead
\bottomrule\noalign{}
\endlastfoot
Finite Weyl equivalence and scalar commutant & Primitive root, unitary
\(N \times N\) pair, period and normalized eigenvalue-one clock seed;
explicit eigenvector-orbit proof in Target 19 & Written theorem and
exact verifier; Lean checks the stated commutant fragment \\
Born normalization & Supplied positive trace-one state, POVM summing to
identity and trace pairing; linearity of trace & General algebraic
identity; specialized two-level Lean models and exact channel checks \\
Gleason\allowbreak{}/\allowbreak{}Busch representation {[}1,3{]} & Full
projective domain at dimension at least three, or an all-effects
additive positive normalized domain & Imported theorems; rational-plane
obstruction and finite tomography do not prove a general number-field
representation theorem \\
Phase-blind reports can lose future answers & Two equal present reports
separated by an admitted common future experiment & Explicit exact
counterexample; no exclusion of every classical enlarged model \\
Bell joint-extension obstruction & Four supplied pair tables and one
positive joint law on all sixteen local assignments & Finite convex
separation and CHSH proof; a physical Bell reading also assumes setting
independence and local response factorization \\
Classical-storage survival at most \(3/4\) & Full qubit positive
states/effects, affine unconditional encoding/retrieval, unknown input
and no quantum side memory & General written proof with sharp control in
§13.7; exact native instances \\
Central-sector purification obstruction and bounds & Full states of
direct-sum complex matrix algebras and standard tensor composition &
General written proof in §13.8; sharpness needs sufficient ancilla
capacity; Lean assumes the pure-block representation \\
Native history quotient & A declared finite reachable
history\allowbreak{}/\allowbreak{}action system and protected future
questions, including availability & Minimal report relative to those
questions; no quantum scalar, positivity or composition selection \\
\end{longtable}

Repository-only checks validate status vocabulary, evidence paths and
named Rust test functions. They enforce bookkeeping; they do not decide
whether a file establishes a claim. The source-by-source discussion in
the source-notes supplement (its §4, subsections 11.1--11.60) supplies
that interpretation and retains counterexamples and open obligations.
Its old execution counts are historical unless an explicitly dated
current replay is named.

\subsection{3. Claims and supporting
evidence}\label{claims-and-supporting-evidence}

The evidence obligation is determined by the statement's quantifiers. A
general mathematical statement requires an adequate written proof or a
formal proof of the same statement. A claim over a finite domain may
instead be settled by a complete exact enumeration, provided the domain
and enumeration are justified. An example test establishes only the
stated example.

Accordingly, each unresolved statement must receive one of the following
treatments:

\begin{enumerate}
\def\labelenumi{\arabic{enumi}.}
\tightlist
\item
  an adequate written proof;
\item
  a matching formal proof;
\item
  a complete exact finite argument appropriate to its quantifiers;
\item
  a restriction to the domain actually established;
\item
  an explicit statement that it remains open; or
\item
  withdrawal after a counterexample.
\end{enumerate}

File existence, compilation, schema validation, serialization, access
control, and a change of status label cannot decide whether this
obligation has been discharged. The repository maintains a historical
evidence register for reproduction, but its row counts and labels are
not a measure of mathematical progress and are not part of this
monograph's scientific summary.

The sixty-seven source manuscripts are preserved in the paper archive.
Six are redirect, alias or annex stubs whose mathematical content is
treated with the destination result: the residual-density redirect
(source-notes supplement §11.4), the finite-realization redirect
(§11.6), the earlier phase-selector alias (§11.27), and the annex stubs
for phase-gauge residuals (§§11.28--11.29), the no-signaling model
(§§11.33--11.34) and the completely positive import boundary (§11.38).

\subsection{4. Established results}\label{established-results}

The results below are stated by mathematical content rather than by
repository status. Their detailed hypotheses, arguments, and source
addresses appear in §§11--13.

\begin{enumerate}
\def\labelenumi{\arabic{enumi}.}
\tightlist
\item
  \textbf{Finite Weyl equivalence from an anchored seed (§1.3, Target
  19).} Given the stated conjugation-stable scalar field
  \(K \subseteq \mathbb C\), primitive \(N\)th root \(\omega\), unitary
  \(N \times N\) pair \(X', Z'\) with \((X')^N = I\), and normalized
  clock eigenvector \(Z'v = v\) with \(v^\dagger v = 1\), the orbit
  columns form an orthonormal eigenbasis and give an explicit unitary
  intertwiner \(W\). The scalar-commutant conclusion on \(K^N\) follows
  by the displayed basis argument. Finding or normalizing such a seed
  over an arbitrary number field is not included; on a larger carrier
  the orbit spans an invariant cyclic subspace.
\item
  \textbf{Finite cochain certificates (source-notes supplement §§11.3,
  11.4 and 11.8).} An exact potential establishes membership in
  \(\operatorname{im}d_0\); a separating functional vanishing on that
  image establishes non-exactness; and a nonzero \(d_1x\) establishes
  failure of closure. Numerical linear programming is reported
  separately.
\item
  \textbf{Power-family context dependence (source-notes supplement
  §11.58, positive-power context independence).} Two rational contexts
  assign a shared ray the weights \(1/2\) and \(1/(1+(3/5)^p+(4/5)^p)\).
  For positive real \(p\), equality holds exactly at \(p=2\). This
  classifies the specified family, not all frame functions.
\item
  \textbf{Positive Bell extension (§13.8).} The full-support rational
  table has CHSH value \(14/5\). Every deterministic assignment is
  support-compatible, yet no nonnegative common joint distribution
  reproduces the four marginals. The separating inequality is verified
  on all sixteen deterministic vertices.
\item
  \textbf{Classical-storage survival bound (§13.7).} Under the stated
  full qubit state\allowbreak{}/\allowbreak{}effect domain, affine
  unconditional measure-and-prepare transfer, unknown input, and absence
  of quantum side memory, the average survival on four preparations is
  at most \(3/4\), and the bound is attained.
\item
  \textbf{Central-sector bounds (§13.8).} In the specified direct-sum
  matrix-algebra state space with standard tensor composition, central
  weights bound extension purity and the trace distance from a
  pure-composite marginal. Attainability requires sufficient ancilla
  capacity.
\item
  \textbf{Primal and dual trace-norm certificates (source-notes
  supplement §11.35).} For \(H=P-N\), positivity of \(P,N\), the
  inequalities \(-I\preceq W\preceq I\), and
  \[\protect\adjustbox{max width=\linewidth}{\(\displaystyle \operatorname{tr}(WH)=\operatorname{tr}P+\operatorname{tr}N\)}\]
  certify the claimed value independently of the implementation's label.
\item
  \textbf{Phase-blind descent with a retained zero point (§11.1).} When
  the quotient keeps the zero vector as a distinct output, the
  population report of \(Ac\) depends only on that of \(c\) exactly when
  every row of \(A\) has at most one nonzero entry. On projective space
  alone the rank-one all-ones matrix is a counterexample. Written proof
  and exact two-dimensional replay.
\item
  \textbf{Trace-zero repair duality (§11.2).} The dual norm of
  \(\frac12\|\cdot\|_1\) on \(\operatorname{Herm}_0(n)\) is the spectral
  diameter. On \(K_{2,3}\) with the stated qutrit defects the trace-zero
  repair value is \(6\) and the all-Hermitian value is \(9/2\), each
  certified by an exact primal potential and a dual witness {[}13{]}.
\item
  \textbf{CPTP monotonicity selects the trace norm (§11.3).} A family of
  norms on trace-zero Hermitian matrices that no CPTP map increases is
  one positive multiple of the trace norm. Written proof by exact
  measure-and-prepare transport.
\item
  \textbf{The Cabello frame polytope (§11.5).} The frame polytope of the
  eighteen Cabello rays has exactly 146 vertices, all half-integral, and
  none of their unique trace-one real-symmetric representatives is
  positive semidefinite. Exact enumeration with
  characteristic-polynomial certificates.
\end{enumerate}

\subsection{5. Open problems and
limitations}\label{open-problems-and-limitations}

The unresolved scientific questions are stated here directly rather than
by internal tracking codes. A question is closed only by a dated
argument or replayable result; renaming or reclassifying it does not
resolve it.

The seven principal finite-to-continuum boundaries are:

\begin{enumerate}
\def\labelenumi{\arabic{enumi}.}
\tightlist
\item
  \textbf{The continuum Hilbert lift.} No source constructs the passage
  from a finite presented carrier to an infinite-dimensional Hilbert
  space. Source Notes §11.52 explicitly refuses that conclusion.
\item
  \textbf{Born uniqueness (Gleason {[}1{]}).} The Born rule is imported
  as a declared input throughout, never derived.
\item
  \textbf{Continuous time and Stone's theorem.} The finite examples
  distinguish finite-order integral evolutions from specified
  infinite-order rational rotations. Rational unitaries can also have
  finite order (the identity and rational quarter-turn suffice). A dense
  cyclic orbit has a closure only after choosing an ambient topology;
  that closure neither selects a time parameter nor a unique logarithm.
  Topology, continuous interpolation and its generator therefore remain
  additional inputs in both sectors.
\item
  \textbf{The Gleason dimension condition at \(n \ge 3\).} Gleason's
  theorem {[}1{]} uses every orthonormal basis of a space of dimension
  at least three, and Busch's version {[}3{]} uses every effect. A
  finite ray set gives only part of that conclusion. On the Cabello set
  in \(\mathbb Q^4\), on its three- and six-basis prefixes, and on the
  24-ray Peres set, every function that sums to one on each declared
  basis is the restriction of a real-symmetric trace-one form (exact
  rank computation). The form is unique once the ray projectors span the
  ten-dimensional space of real symmetric \(4\times4\) matrices, which
  holds on every listed set except the three-basis prefix. Nonnegativity
  on the declared rays does not give positive semidefiniteness: the 146
  vertex representatives of §11.5 are nonnegative on the eighteen
  Cabello rays, the research note
  \texttt{r\allowbreak{}e\allowbreak{}s\allowbreak{}e\allowbreak{}a\allowbreak{}r\allowbreak{}c\allowbreak{}h\allowbreak{}/\allowbreak{}q\allowbreak{}m\allowbreak{}\_\allowbreak{}2\allowbreak{}1\allowbreak{}0\allowbreak{}\_\allowbreak{}f\allowbreak{}r\allowbreak{}a\allowbreak{}m\allowbreak{}e\allowbreak{}\_\allowbreak{}f\allowbreak{}u\allowbreak{}n\allowbreak{}c\allowbreak{}t\allowbreak{}i\allowbreak{}o\allowbreak{}n\allowbreak{}\_\allowbreak{}f\allowbreak{}i\allowbreak{}n\allowbreak{}i\allowbreak{}t\allowbreak{}e\allowbreak{}\_\allowbreak{}m\allowbreak{}a\allowbreak{}p\allowbreak{}.\allowbreak{}m\allowbreak{}d}
  displays a rational trace-one form nonnegative on both the Cabello and
  the Peres rays, and none of these is positive semidefinite. The finite
  data fix the linear representative but not the positive cone, so the
  dimension condition remains an imported premise.
\item
  \textbf{The rational boundary of Everett-style branch counting.}
  Branch counting on fine-grained trees reproduces Born weights if and
  only if all weights lie in \(\mathbb{Q}\). For weights in
  \[\protect\adjustbox{max width=\linewidth}{\(\displaystyle \mathbb{Q}(\zeta_n)^+ \setminus \mathbb{Q}\)}\],
  finite integer fine-graining encounters an exact arithmetic
  obstruction.
\item
  \textbf{Choice of real embedding.} The maximal totally real subfield
  \(\mathbb{Q}(\zeta_n)^+\) has \(\phi(n)/2\) real embeddings into
  \(\mathbb{R}\) (16 at target conductor \(n = 120\)). Gram matrices of
  pure states stay positive semidefinite under every Galois conjugation,
  since each conjugate is again a Gram matrix, but mixed-state witnesses
  (starting at \(n = 5\)) exhibit embedding-relative positivity. The
  Galois group acts transitively on embeddings; selecting an ordered
  real embedding is an external declaration.
\item
  \textbf{The transcendental von Neumann entropy interface for mixed
  states.} For a non-pure density operator with algebraic spectrum,
  entropy in natural-log units,
  \[\protect\adjustbox{max width=\linewidth}{\(\displaystyle S=-\sum_{\lambda_i>0}\lambda_i\ln\lambda_i\)}\],
  is positive and is a nonzero algebraic linear combination of
  logarithms of algebraic numbers. Baker's theorem on linear forms in
  logarithms {[}24,27{]} makes this value transcendental. The claim is
  unit-sensitive: in bits, \(S(I_2/2)=1\), so it does not assert that
  every base-two entropy is transcendental. The natural-log quantity
  lies outside \(\mathbb Q(\zeta_n)\); the declared external interface
  uses a certified rational enclosure \([L_k,U_k]\) of width at most
  \(2^{-k}\), with its midpoint error bounded by half that width.
\end{enumerate}

If finite quantum mechanics fails as a reconstruction programme, the
failure is most likely to be exhibited at one of these seven boundaries.
The useful outcome would be a precise theorem locating that boundary,
not a change in a project status label. Historical packaging corrections
and evidence-register statistics are retained in the source notes and
reproduction records. They are not measures of the mathematical content
established here.

\subsubsection{5.1 Executable examples do not close the derivation
gap}\label{executable-examples-do-not-close-the-derivation-gap}

The repository contains executable interferometer, finite-cochain,
Cayley-step, and positive-semidefinite filtering examples. They consume
a declared graph, scalar field, state or amplitude data, readout rule,
and experiment. Their successful execution establishes facts about those
inputs; it does not derive the inputs or their quantum interpretation
from primitive difference.

The exact examples retained in the main text use rational or specified
algebraic arithmetic. Older numerical linear-programming examples remain
numerical demonstrations. Historical execution counts and migration
details are recorded in the source-notes supplement, not used as theorem
counts.

\subsection{6. How to read the publication
package}\label{how-to-read-the-publication-package}

The main paper states current mathematical results and limitations. The
source-notes supplement preserves the paper-by-paper corrections,
original source addresses, and detailed implementation scope. All
retained source manuscripts cited there are included under
\texttt{s\allowbreak{}o\allowbreak{}u\allowbreak{}r\allowbreak{}c\allowbreak{}e\allowbreak{}s\allowbreak{}/\allowbreak{}o\allowbreak{}r\allowbreak{}i\allowbreak{}g\allowbreak{}i\allowbreak{}n\allowbreak{}a\allowbreak{}l\allowbreak{}\_\allowbreak{}p\allowbreak{}a\allowbreak{}p\allowbreak{}e\allowbreak{}r\allowbreak{}s\allowbreak{}/}
in the paper ZIP. Companion papers used in the current exposition are
included under
\texttt{c\allowbreak{}o\allowbreak{}m\allowbreak{}p\allowbreak{}a\allowbreak{}n\allowbreak{}i\allowbreak{}o\allowbreak{}n\allowbreak{}\_\allowbreak{}p\allowbreak{}a\allowbreak{}p\allowbreak{}e\allowbreak{}r\allowbreak{}s\allowbreak{}/}.

The separate code-companion ZIP contains the shared Rust and Lean source
trees in the same relative layout used by their manifests. This split
keeps the scientific paper readable while allowing a reader to inspect
and run the actual code instead of relying on absent repository paths.
Section 12 gives the focused commands and distinguishes mathematical
proof, finite replay, formal checking, and archive custody.

For an unresolved statement, the obligation is unchanged: provide an
adequate written proof, a formal proof of the same statement, or a
complete exact finite argument appropriate to its quantifiers; otherwise
restrict it, identify it as open, or withdraw it. An example test
remains evidence only for its stated inputs.

\subsection{7. Scope limits}\label{scope-limits}

\begin{itemize}
\tightlist
\item
  Nothing here derives the Born rule, complex structure, unitarity,
  tensor factorization, or continuous time. Every one of those enters
  the source papers as a \textbf{declared import}, and importing a thing
  is not deriving it.
\item
  A green test suite in this application certifies that its mechanical
  constraints hold --- that files exist, commands exit 0, syntax parses,
  and unit tests execute. It does not by itself certify semantic
  mathematical validity or that prose claims are fully entailed by
  code\allowbreak{}/\allowbreak{}proof terms; those boundaries require
  substantive review.
\end{itemize}

\subsection{8. Sources, code, and public-record
corrections}\label{sources-code-and-public-record-corrections}

The companion publication ZIP contains the retained source manuscripts
under
\texttt{s\allowbreak{}o\allowbreak{}u\allowbreak{}r\allowbreak{}c\allowbreak{}e\allowbreak{}s\allowbreak{}/\allowbreak{}o\allowbreak{}r\allowbreak{}i\allowbreak{}g\allowbreak{}i\allowbreak{}n\allowbreak{}a\allowbreak{}l\allowbreak{}\_\allowbreak{}p\allowbreak{}a\allowbreak{}p\allowbreak{}e\allowbreak{}r\allowbreak{}s\allowbreak{}/},
the source-notes supplement, and the code used by the package-local
reproductions. A separate code-companion ZIP preserves the shared Rust
workspace and the Lean workspace in their repository-relative locations
so the named commands do not depend on files absent from the download.
Section 12 states the exact commands and the limits of the latest
replay.

Two earlier Zenodo records require author-controlled correction before
this edition is promoted. Prepared errata are included locally:

\begin{itemize}
\tightlist
\item
  \texttt{d\allowbreak{}o\allowbreak{}c\allowbreak{}s\allowbreak{}/\allowbreak{}z\allowbreak{}e\allowbreak{}n\allowbreak{}o\allowbreak{}d\allowbreak{}o\allowbreak{}\_\allowbreak{}e\allowbreak{}r\allowbreak{}r\allowbreak{}a\allowbreak{}t\allowbreak{}a\allowbreak{}/\allowbreak{}2\allowbreak{}0\allowbreak{}6\allowbreak{}8\allowbreak{}3\allowbreak{}7\allowbreak{}7\allowbreak{}5\allowbreak{}\_\allowbreak{}e\allowbreak{}r\allowbreak{}r\allowbreak{}a\allowbreak{}t\allowbreak{}u\allowbreak{}m\allowbreak{}.\allowbreak{}m\allowbreak{}d}
  narrows the claimed relationship between phase averaging, the Born
  rule, and Hilbert-space emergence;
\item
  \texttt{d\allowbreak{}o\allowbreak{}c\allowbreak{}s\allowbreak{}/\allowbreak{}z\allowbreak{}e\allowbreak{}n\allowbreak{}o\allowbreak{}d\allowbreak{}o\allowbreak{}\_\allowbreak{}e\allowbreak{}r\allowbreak{}r\allowbreak{}a\allowbreak{}t\allowbreak{}a\allowbreak{}/\allowbreak{}2\allowbreak{}0\allowbreak{}5\allowbreak{}5\allowbreak{}0\allowbreak{}4\allowbreak{}5\allowbreak{}9\allowbreak{}\_\allowbreak{}e\allowbreak{}r\allowbreak{}r\allowbreak{}a\allowbreak{}t\allowbreak{}u\allowbreak{}m\allowbreak{}.\allowbreak{}m\allowbreak{}d}
  replaces the incorrect centered trace-zero dual and records the exact
  \(K_{2,3}\) counterexample.
\end{itemize}

These files are proposed public-record corrections, not evidence that
Zenodo has already been updated. The source-notes supplement contains
the full version and custody discussion.

\subsection{9. Limits of the proposed finite
derivation}\label{limits-of-the-proposed-finite-derivation}

An earlier version of this section claimed a complete 50-step
derivation, a minimal ambient carrier, and resolution of every proposed
test. The proofs and computations do not establish those claims, so they
are withdrawn. The historical wording remains available in version
control.

\begin{longtable}[]{@{}
  >{\raggedright\arraybackslash}p{(\linewidth - 4\tabcolsep) * \real{0.3333}}
  >{\raggedright\arraybackslash}p{(\linewidth - 4\tabcolsep) * \real{0.3333}}
  >{\raggedright\arraybackslash}p{(\linewidth - 4\tabcolsep) * \real{0.3333}}@{}}
\toprule\noalign{}
\begin{minipage}[b]{\linewidth}\raggedright
Topic
\end{minipage} & \begin{minipage}[b]{\linewidth}\raggedright
What the Lean statements cover
\end{minipage} & \begin{minipage}[b]{\linewidth}\raggedright
What they do not cover
\end{minipage} \\
\midrule\noalign{}
\endhead
\bottomrule\noalign{}
\endlastfoot
Carrier labels & Natural-number divisibility, a directed label system,
and Gaussian record identities & An actual field colimit, all-target
embedding, or minimal carrier \\
Probability and contextuality & Declared squared-norm weights, two
Sorkin fixtures, and finite CHSH\allowbreak{}/\allowbreak{}parity
arithmetic & Quotient-fiber transport, marginal no-signaling, or a
derivation of Born weights \\
Cabello witness & Cabello matrix checks and a norm-25 integer partition
& Topology or non-Bornian functions on the dense cyclotomic sphere \\
Discrete dynamics & Boolean endomaps, a scaled \(2\times2\) Gaussian
square root, integer Cayley identities, and integer-labelled records &
Continuous Hamiltonians, logarithms, Stone's theorem, or a necessity
proof \\
Commutator obstruction & A \(2\times2\) integer commutator trace
obstruction and one Weyl example & All dimensions and rings, a
completion functor, or an exclusive continuum classification \\
Conductor bounds & Divisibility and least common multiples of declared
labels & State embeddings, field minimality, or a general stabilizer or
unextendible-product-basis theorem \\
\end{longtable}

The Python Cabello witness satisfies the
18-ray\allowbreak{}/\allowbreak{}9-context constraints and has a checked
negative eigenvector. The Peres computation checks 33 rays, 72 pairs and
16 triads; its binary valuation count is zero. Its finite linear system
has rank 16. For the constructed function
\texttt{f\allowbreak{}=\allowbreak{}1\allowbreak{}/\allowbreak{}3\allowbreak{}+\allowbreak{}z},
the residual \texttt{z}, not \texttt{f}, is orthogonal to the declared
trace-free quadratic rows. Neither construction is a frame function on a
dense sphere.

Python constructs ten bounded triads. Rust's legacy higher-conductor
report checks four rational triads only; its former six additional
offers were not computed and are withdrawn. No cross-language
higher-triad parity is claimed. The topology, rational-sphere
colorability, p-adic and nonstandard-model flags that lacked
implementations are now explicitly unknown. Modular witnesses and
four-conductor discriminant checks retain their finite scope.

Preserving an earlier prediction text does not validate its conclusion.
The proposed refusals are not established, and both Gleason-related
questions remain open under explicit conditions. No ``probability wall
is absolute'' or ``completion is necessary'' conclusion survives this
review.

The strongest result here is a collection of exact finite computations
and scoped Lean theorems under declared carriers and measures.
Finite-dimensional does not mean finite carrier; countable does not mean
discrete topology.

\subsection{10. Higher-dimensional bounds and the Tiles
construction}\label{higher-dimensional-bounds-and-the-tiles-construction}

Conductor labels here are sufficient declarations for the named objects,
not a search proving minimality. In particular, a generator or
parity-check matrix is not the same object as a normalized state.

\begin{itemize}
\tightlist
\item
  The five-qubit stabilizer
  generator\allowbreak{}/\allowbreak{}projector declaration uses
  \(\mathbb Q(i)\), label 4; the Steane parity-check declaration uses
  \(\mathbb Q\), label 1. No normalized-codeword carrier or
  all-stabilizer classification follows.
\item
  The named GHZ and W normalization declarations use labels 8 and 12. A
  seven-qubit W normalization uses the positive root \(\sqrt7\),
  discriminant 28.
\item
  Prime-qudit normalized setting declarations combine roots of unity and
  the positive root \(\sqrt p\). At \(p=7\) the bound uses 28 and
  \(\operatorname{lcm}(120,28)=840\). (The field
  \(\mathbb Q(\sqrt{-7})\) has discriminant \(-7\), not 28.)
\item
  Collins--Gisin--Linden--Massar--Popescu (CGLMP) setting bounds use
  \[\protect\adjustbox{max width=\linewidth}{\(\displaystyle \operatorname{lcm}(4d,\operatorname{disc}\mathbb Q(\sqrt d))\)}\];
  they exclude an optimal entangled state's amplitudes. For
  \(d=3,4,5,6,7,8\) the labels are respectively 12, 16, 20, 24, 28, 32,
  giving joint least common multiples with 120 of 120, 240, 120, 120,
  840, 480.
\end{itemize}

The Tiles \(3\times3\) example is actual matrix data, not a trace-only
assertion. With five unnormalized product vectors \(v_i\), it constructs
the rational matrix

\[\protect\adjustbox{max width=\linewidth}{\(\displaystyle \rho=\frac14\left(I_9-\sum_i\frac{v_iv_i^{\mathsf T}}{v_i^{\mathsf T}v_i}\right).\)}\]

Python and Rust verify mutual orthogonality, trace one, symmetry, the
projector identity
\texttt{4\allowbreak{} \allowbreak{}r\allowbreak{}h\allowbreak{}o\allowbreak{}\textasciicircum{}\allowbreak{}2\allowbreak{} \allowbreak{}=\allowbreak{} \allowbreak{}r\allowbreak{}h\allowbreak{}o},
and equality with its partial transpose. All 32 assignments of local
orthogonality conditions are checked by exact rank-three minors,
excluding a nonzero product vector in the complement. The entanglement
reading uses the declared UPB-complement criterion; this is a
deterministic finite certificate, not a Lean UPB theorem or general
classifier. Duplicate and zero product-vector inputs are rejected.

From
\texttt{p\allowbreak{}y\allowbreak{}t\allowbreak{}h\allowbreak{}o\allowbreak{}n\allowbreak{}/},
\texttt{p\allowbreak{}y\allowbreak{}t\allowbreak{}h\allowbreak{}o\allowbreak{}n\allowbreak{} \allowbreak{}-\allowbreak{}m\allowbreak{} \allowbreak{}f\allowbreak{}i\allowbreak{}n\allowbreak{}i\allowbreak{}t\allowbreak{}e\allowbreak{}\_\allowbreak{}q\allowbreak{}m\allowbreak{}.\allowbreak{}e\allowbreak{}n\allowbreak{}t\allowbreak{}r\allowbreak{}y\allowbreak{}p\allowbreak{}o\allowbreak{}i\allowbreak{}n\allowbreak{}t\allowbreak{}s\allowbreak{}.\allowbreak{}r\allowbreak{}e\allowbreak{}p\allowbreak{}l\allowbreak{}a\allowbreak{}y\allowbreak{}\_\allowbreak{}r\allowbreak{}e\allowbreak{}v\allowbreak{}i\allowbreak{}e\allowbreak{}w\allowbreak{}\_\allowbreak{}r\allowbreak{}e\allowbreak{}c\allowbreak{}e\allowbreak{}i\allowbreak{}p\allowbreak{}t\allowbreak{}s\allowbreak{} \allowbreak{}-\allowbreak{}-\allowbreak{}c\allowbreak{}h\allowbreak{}e\allowbreak{}c\allowbreak{}k}
regenerates both repaired computation receipts byte-for-byte. Receipt
identity does not promote the declarations above.

\subsection{11. Principal exact results}\label{principal-exact-results}

This section collects the strongest results that emerged from the exact
finite analysis. The long source-by-source record is in
\texttt{f\allowbreak{}i\allowbreak{}n\allowbreak{}i\allowbreak{}t\allowbreak{}e\allowbreak{}\_\allowbreak{}q\allowbreak{}u\allowbreak{}a\allowbreak{}n\allowbreak{}t\allowbreak{}u\allowbreak{}m\allowbreak{}\_\allowbreak{}m\allowbreak{}e\allowbreak{}c\allowbreak{}h\allowbreak{}a\allowbreak{}n\allowbreak{}i\allowbreak{}c\allowbreak{}s\allowbreak{}\_\allowbreak{}s\allowbreak{}o\allowbreak{}u\allowbreak{}r\allowbreak{}c\allowbreak{}e\allowbreak{}\_\allowbreak{}n\allowbreak{}o\allowbreak{}t\allowbreak{}e\allowbreak{}s\allowbreak{}.\allowbreak{}m\allowbreak{}d}
(\texttt{f\allowbreak{}i\allowbreak{}n\allowbreak{}i\allowbreak{}t\allowbreak{}e\allowbreak{}\_\allowbreak{}q\allowbreak{}u\allowbreak{}a\allowbreak{}n\allowbreak{}t\allowbreak{}u\allowbreak{}m\allowbreak{}\_\allowbreak{}m\allowbreak{}e\allowbreak{}c\allowbreak{}h\allowbreak{}a\allowbreak{}n\allowbreak{}i\allowbreak{}c\allowbreak{}s\allowbreak{}\_\allowbreak{}s\allowbreak{}o\allowbreak{}u\allowbreak{}r\allowbreak{}c\allowbreak{}e\allowbreak{}\_\allowbreak{}n\allowbreak{}o\allowbreak{}t\allowbreak{}e\allowbreak{}s\allowbreak{}.\allowbreak{}m\allowbreak{}d}).
Each result below is stated by its mathematical content. Implementation
names are supplied only as reproduction pointers.

\subsubsection{11.1 Phase-blind descent requires a retained zero
point}\label{phase-blind-descent-requires-a-retained-zero-point}

Let \(A\) be a matrix acting on nonzero vectors and let

\[\protect\adjustbox{max width=\linewidth}{\(\displaystyle \pi(c)_g=\frac{|c_g|^2}{\sum_h|c_h|^2}\)}\]

be the phase-blind population report. A tempting claim says that
\(\pi(Ac)\) depends only on \(\pi(c)\) exactly when every row of \(A\)
has at most one nonzero entry. That claim is false on projective space.
A rank-one matrix such as

\[\protect\adjustbox{max width=\linewidth}{\(\displaystyle A=\begin{pmatrix}1&1\\1&1\end{pmatrix}\)}\]

sends every non-annihilated input to the same ray even though each row
has two nonzero entries.

The row-sparsity characterization becomes correct when the quotient
retains the zero vector as a distinct output. If two nonzero entries
occur in one row, relative phase can be chosen so that this row cancels
for one representative but not for another representative with the same
population report. The retained zero point detects that annihilation.
Conversely, a row-sparse matrix cannot create such phase-sensitive cross
terms. The exact two-dimensional counterexamples and the zero-retaining
proof are replayed in
\texttt{p\allowbreak{}h\allowbreak{}a\allowbreak{}s\allowbreak{}e\allowbreak{}\_\allowbreak{}b\allowbreak{}l\allowbreak{}i\allowbreak{}n\allowbreak{}d\allowbreak{}\_\allowbreak{}q\allowbreak{}u\allowbreak{}o\allowbreak{}t\allowbreak{}i\allowbreak{}e\allowbreak{}n\allowbreak{}t}
and described in source-notes supplement §11.25.

\subsubsection{11.2 Trace-zero repair quotients have a scalar-divergence
dual}\label{trace-zero-repair-quotients-have-a-scalar-divergence-dual}

For Hermitian trace-zero edge defects, compare

\[\protect\adjustbox{max width=\linewidth}{\(\displaystyle \Phi_0(\delta)=\inf_{\alpha_v\in\operatorname{Herm}_0(n)}
\sum_e\frac12\|\delta_e-(d_0\alpha)_e\|_1\)}\]

with the same infimum over all Hermitian vertex potentials. The polar of
\(\frac12\|\cdot\|_1\) on \(\operatorname{Herm}_0(n)\) is an observable
class modulo scalar identity, with norm

\[\protect\adjustbox{max width=\linewidth}{\(\displaystyle \|[M]\|^*=\lambda_{\max}(M)-\lambda_{\min}(M).\)}\]

Accordingly, the trace-zero dual requires spectral diameter at most one
and \((d_0^*M)_v\in\mathbb RI\). The all-Hermitian dual instead requires
\(\|M_e\|_{\mathrm{op}}\leq1/2\) and \(d_0^*M=0\).

On \(K_{2,3}\), orient the legs \(u\to w_k\to v\) and put
\(D_k=3E_{kk}-I\) on \(u\to w_k\). Zero repair and the witnesses
\(D_k/3\) on both edges of each leg certify

\[\protect\adjustbox{max width=\linewidth}{\(\displaystyle \Phi_0=6.\)}\]

The all-Hermitian potential \(\alpha_u=0\), \(\alpha_v=-I\),
\(\alpha_{w_k}=D_k\), together with the witnesses \(D_k/4\), certifies

\[\protect\adjustbox{max width=\linewidth}{\(\displaystyle \Phi_{\mathrm{all}}=\frac92.\)}\]

The complete theorem and dependency-free exact replay are in the
companion paper \emph{Trace-Zero Repair Quotients on a Finite Graph}
{[}13{]}.

\subsubsection{11.3 CPTP monotonicity selects the trace norm on state
differences}\label{cptp-monotonicity-selects-the-trace-norm-on-state-differences}

For each finite dimension \(n\), let \(N_n\) be a norm on the real
vector space
\[\protect\adjustbox{max width=\linewidth}{\(\displaystyle \operatorname{Herm}_0(n, \mathbb C)\)}\].
Suppose

\[\protect\adjustbox{max width=\linewidth}{\(\displaystyle N_m(\Phi(X)) \le N_n(X)\)}\]

for every CPTP map \(\Phi\) from \(M_n(\mathbb C)\) to
\(M_m(\mathbb C)\) and every
\[\protect\adjustbox{max width=\linewidth}{\(\displaystyle X \in \operatorname{Herm}_0(n, \mathbb C)\)}\].
Then there is a single \(c > 0\) such that \(N_n(X) = c \|X\|_1\) on
every nonzero such space. The dimension-one space is zero and satisfies
the identity trivially.

To prove this, let \(X = X_+ - X_-\) be nonzero and set
\[\protect\adjustbox{max width=\linewidth}{\(\displaystyle t = \operatorname{tr} X_+ = \|X\|_1 / 2 > 0\)}\].
For any target \(Y = Y_+ - Y_-\) with
\(s = \operatorname{tr} Y_+ \le t\), choose an output density matrix
\(\tau\) and put

\[\protect\adjustbox{max width=\linewidth}{\(\displaystyle A = \frac{Y_+}{t} + \left(1 - \frac{s}{t}\right)\tau, \qquad B = \frac{Y_-}{t} + \left(1 - \frac{s}{t}\right)\tau.\)}\]

These are normalized positive states. If \(P\) is the positive spectral
support projection of \(X\), the map

\[\protect\adjustbox{max width=\linewidth}{\(\displaystyle \Phi(T) = \operatorname{tr}(PT)A + \operatorname{tr}((I-P)T)B\)}\]

is CPTP and satisfies \(\Phi(X) = t(A - B) = Y\) exactly. Equal trace
norms admit exact transports in both directions, hence the \(N\)-values
are equal. Absolute homogeneity then gives the common constant \(c\).
This is a theorem on the full stated channel and matrix domain, not a
selection of that domain from primitive difference or an assertion that
every construction stays inside an arbitrary smaller scalar field. Exact
difference-transport examples are in source-notes supplement §11.36.

\subsubsection{11.4 Full-support Bell tables require positivity, not
support
alone}\label{full-support-bell-tables-require-positivity-not-support-alone}

Four rational pair distributions can have full support while admitting
no positive common joint law. The supplied qubit construction uses a
rational density matrix and \((3,4,5)\)-triangle observables and has

\[\protect\adjustbox{max width=\linewidth}{\(\displaystyle S_{\mathrm{CHSH}}=\frac{14}{5}>2.\)}\]

Every deterministic assignment is compatible with the support, so a
possibilistic classifier returns no obstruction. A positive common
distribution would be a convex combination of the sixteen deterministic
vertices, each of which satisfies \(|S|\leq2\); therefore no such
distribution exists. This cleanly separates support gluing from positive
probabilistic extension. The broader contextuality and no-signaling
scope is in §13.8.

\subsubsection{11.5 The Cabello frame polytope has 146 half-integral
vertices}\label{the-cabello-frame-polytope-has-146-half-integral-vertices}

Let \(v_1,\ldots,v_{18}\) be the eighteen integer rays

\[\protect\adjustbox{max width=\linewidth}{\(\displaystyle \begin{gathered}
(1,1,1,1),\ (1,1,1,-1),\ (1,1,0,0),\ (1,1,-1,-1),\ (1,0,1,0),\ (1,0,0,1),\\
(1,0,0,-1),\ (1,0,-1,0),\ (1,-1,1,1),\ (1,-1,0,0),\ (1,-1,-1,1),\ (1,-1,-1,-1),\\
(0,1,0,1),\ (0,1,0,0),\ (0,1,-1,0),\ (0,0,1,0),\ (0,0,1,-1),\ (0,0,0,1),
\end{gathered}\)}\]

and let the nine tetrads be their orthogonal four-element subsets. Every
ray lies in exactly two tetrads, and the orthogonality graph and tetrad
incidence are isomorphic to those of the Cabello configuration {[}2{]}.
Put \(P_i = v_i v_i^T / (v_i^T v_i)\), and let \(A\) be its
\(9 \times 18\) tetrad incidence matrix. The frame polytope is
\[\protect\adjustbox{max width=\linewidth}{\(\displaystyle \mathcal P = \{f \in \mathbb R^{18} : f \ge 0,\ Af = \mathbf 1\}\)}\].
Its real-symmetric operator representative \(H_f\) is defined by
\(\operatorname{tr}(H_f P_i) = f_i\). The ten-dimensional real-symmetric
projector span and the rank-nine tetrad constraints give the unique
trace-one representative.

Every extreme valuation is obtained by covering the nine tetrad vertices
with disjoint isolated edges of weight one and odd cycles of weight one
half. An exact enumeration yields 146 valuations. Their characteristic
polynomials and multiplicities are:

\begin{longtable}[]{@{}
  >{\raggedright\arraybackslash}p{(\linewidth - 4\tabcolsep) * \real{0.3333}}
  >{\raggedright\arraybackslash}p{(\linewidth - 4\tabcolsep) * \real{0.3333}}
  >{\raggedright\arraybackslash}p{(\linewidth - 4\tabcolsep) * \real{0.3333}}@{}}
\toprule\noalign{}
\begin{minipage}[b]{\linewidth}\raggedright
Multiplicity
\end{minipage} & \begin{minipage}[b]{\linewidth}\raggedright
Characteristic polynomial of \(H_f\)
\end{minipage} & \begin{minipage}[b]{\linewidth}\raggedright
Smallest eigenvalue
\end{minipage} \\
\midrule\noalign{}
\endhead
\bottomrule\noalign{}
\endlastfoot
24 & \(t^2(t-3/2)(t+1/2)\) & \(-1/2\) \\
72 & \(t^2(t^2-t-1/4)\) & \((1-\sqrt2)/2\) \\
2 & \((t-1)^2(t+1/2)^2\) & \(-1/2\) \\
48 & \((t^2-t/2-1/8)^2\) & \((1-\sqrt3)/4\) \\
\end{longtable}

Hence no vertex representative is positive semidefinite; 26 have
smallest eigenvalue \(-1/2\), and the other 120 have smallest eigenvalue
greater than \(-1/4\). The complete input, enumeration argument, vertex
list and exact reconstruction checks accompany the result. These
statements concern the real-symmetric representatives, not arbitrary
complex-Hermitian matrices that agree only on the sampled real rays. The
enumeration, rank checks, and exact characteristic-polynomial
certificates are described in source-notes supplement §11.52.

\subsubsection{\texorpdfstring{11.6 A single rational context pair
selects \(p=2\) among power
rules}{11.6 A single rational context pair selects p=2 among power rules}}\label{a-single-rational-context-pair-selects-p2-among-power-rules}

Consider the positive-power family

\[\protect\adjustbox{max width=\linewidth}{\(\displaystyle P_p(i)=\frac{|c_i|^p}{\sum_j|c_j|^p},\qquad p>0.\)}\]

For \(\psi=(1,1,0)\), compare the standard basis with the rational
orthonormal basis obtained by a \((3,4,5)\) rotation fixing the first
coordinate. The shared ray has weight \(1/2\) in the first context and

\[\protect\adjustbox{max width=\linewidth}{\(\displaystyle \frac{1}{1+(3/5)^p+(4/5)^p}\)}\]

in the second. Agreement is equivalent to \((3/5)^p+(4/5)^p=1\). The
left side is strictly decreasing for \(p>0\) and equals one at \(p=2\),
so \(p=2\) is the unique positive solution. This classifies the declared
power family; it is not a general Gleason theorem. Sorkin's third-order
sum {[}7{]} does not make this selection: for the polynomial measures
\((a+b+c)^p\) it vanishes identically at \(p=1\), the additive case, and
at \(p=2\), the quadratic case (source-notes supplement §11.52, check
C46). Exact rational controls for selected powers accompany the general
monotonicity argument in source-notes supplement §11.58.

\subsection{12. Reproduction and package
boundaries}\label{reproduction-and-package-boundaries}

The publication is distributed as two archives with different purposes.

\begin{enumerate}
\def\labelenumi{\arabic{enumi}.}
\tightlist
\item
  The paper archive contains this monograph, the source-notes
  supplement, retained source manuscripts, package-local Python
  reproductions, exact receipts, and checksum manifests.
\item
  The code-companion archive preserves the shared Rust workspace, the
  finite quantum mechanics native crate and campaign crate, and the Lean
  workspace in their repository-relative locations. Its own manifest
  binds every member.
\end{enumerate}

The exact filename, byte count, and SHA-256 digest of the code companion
are recorded in
\texttt{d\allowbreak{}o\allowbreak{}c\allowbreak{}s\allowbreak{}/\allowbreak{}c\allowbreak{}o\allowbreak{}d\allowbreak{}e\allowbreak{}\_\allowbreak{}c\allowbreak{}o\allowbreak{}m\allowbreak{}p\allowbreak{}a\allowbreak{}n\allowbreak{}i\allowbreak{}o\allowbreak{}n\allowbreak{}.\allowbreak{}m\allowbreak{}d}
and bound again by
\texttt{C\allowbreak{}O\allowbreak{}D\allowbreak{}E\allowbreak{}\_\allowbreak{}C\allowbreak{}O\allowbreak{}M\allowbreak{}P\allowbreak{}A\allowbreak{}N\allowbreak{}I\allowbreak{}O\allowbreak{}N\allowbreak{}.\allowbreak{}s\allowbreak{}h\allowbreak{}a\allowbreak{}2\allowbreak{}5\allowbreak{}6}.

From the extracted code-companion root, the focused trace-zero
certificate is run with

\begin{Shaded}
\begin{Highlighting}[]
\NormalTok{cargo test {-}{-}manifest{-}path omega\_engine/crates\_rs/Cargo.toml \textbackslash{}}
\NormalTok{  {-}p omega\_quantum\_physics {-}{-}test channels \textbackslash{}}
\NormalTok{  trace\_zero\_repairs\_need\_uncentred\_witnesses\_on\_two\_cycles\_of\_qutrit\_defects}
\end{Highlighting}
\end{Shaded}

The complete Lean workspace is built with

\begin{Shaded}
\begin{Highlighting}[]
\NormalTok{cd lean4/quantum\_mechanics}
\NormalTok{lake build}
\end{Highlighting}
\end{Shaded}

The dependency-free replay for the standalone \(K_{2,3}\) paper is also
included at
\texttt{c\allowbreak{}o\allowbreak{}m\allowbreak{}p\allowbreak{}a\allowbreak{}n\allowbreak{}i\allowbreak{}o\allowbreak{}n\allowbreak{}\_\allowbreak{}p\allowbreak{}a\allowbreak{}p\allowbreak{}e\allowbreak{}r\allowbreak{}s\allowbreak{}/\allowbreak{}t\allowbreak{}r\allowbreak{}a\allowbreak{}c\allowbreak{}e\allowbreak{}\_\allowbreak{}z\allowbreak{}e\allowbreak{}r\allowbreak{}o\allowbreak{}\_\allowbreak{}q\allowbreak{}u\allowbreak{}o\allowbreak{}t\allowbreak{}i\allowbreak{}e\allowbreak{}n\allowbreak{}t\allowbreak{}\_\allowbreak{}d\allowbreak{}u\allowbreak{}a\allowbreak{}l\allowbreak{}i\allowbreak{}t\allowbreak{}y\allowbreak{}\_\allowbreak{}m\allowbreak{}o\allowbreak{}n\allowbreak{}o\allowbreak{}g\allowbreak{}r\allowbreak{}a\allowbreak{}p\allowbreak{}h\allowbreak{}.\allowbreak{}m\allowbreak{}d};
its package runs with standard-library rational arithmetic and focused
\texttt{p\allowbreak{}y\allowbreak{}t\allowbreak{}e\allowbreak{}s\allowbreak{}t}
tests. The paper archive is checked by

\begin{Shaded}
\begin{Highlighting}[]
\NormalTok{python scripts/regen\_sha256sums.py {-}{-}check}
\NormalTok{python verify\_deposit.py}
\end{Highlighting}
\end{Shaded}

These commands have separate evidence scopes. Exact certificate replay
establishes the displayed finite identities. A Rust or Lean build checks
the statements actually encoded in those sources. Manifest and archive
verification establish byte custody. None of these operations derives
the physical premises named in the manuscript, and this edition does not
claim a fresh successful build of every Rust and Lean target. The full
implementation history and per-module scope are retained in the
source-notes supplement.

\subsection{13. The quantum carrier, observer, and measurement
problem}\label{the-quantum-carrier-observer-and-measurement-problem}

\subsubsection{13.1 The current answer and the questions it leaves
open}\label{the-current-answer-and-the-questions-it-leaves-open}

Our best current answer is a conditional operational account. A quantum
carrier supports the declared states, interactions and effects. An
observer specifies which of those operations are available, how an
apparatus registers a result, and which distinctions its reports must
preserve. Measurement couples carriers and makes information available
in a record. A classical description can then be sufficient for a
specified set of questions even while the larger carrier retains quantum
coherence.

That account separates three objects: \textbf{the probability
distribution of records, the state conditioned on an obtained record,
and one actual event}. Neither a report label nor a reduced density
matrix supplies the last object by itself. A decision rule evaluates a
supplied record; it is not a physical mechanism that makes a particular
outcome occur.

This account does not imply that every observer requires a quantum
carrier. Defining an observer model requires an adequate carrier, but
does not determine complex scalars, a Hilbert space, a tensor product or
the Born rule. Those are currently explicit inputs to the quantum models
in this paper.

There are several measurement problems, with different obligations
{[}9{]}:

\begin{itemize}
\tightlist
\item
  \textbf{Carrier and interaction:} why this state space, composition
  law and measuring dynamics are physically appropriate.
\item
  \textbf{Preferred records:} which alternatives are distinguishable and
  stable under the actual apparatus and environmental couplings.
\item
  \textbf{Probabilities:} why a specified instrument has its stated
  outcome weights, and under which domain and assumptions those weights
  have a Born form.
\item
  \textbf{Definite outcomes:} what accounts for one event rather than
  the joint state or its distribution of possible records.
\item
  \textbf{Agreement:} when records from different observers can be
  compared and composed without assigning incompatible experiments an
  unsupported common history.
\end{itemize}

Decoherence addresses the suppression of accessible interference and the
stability of certain records {[}8, 10{]}. It does not, under unitary
evolution alone, select a single event. Our finite account makes this
boundary explicit and offers tests at it; it does not announce a
solution to every problem above.

\subsubsection{13.2 The carrier and the declared
observer}\label{the-carrier-and-the-declared-observer}

A finite-dimensional quantum model may have continuously many states;
finite dimension, a finite presentation and a finite state set are
different claims. In our exact models, entries lie in a declared field
such as the Gaussian rationals or a cyclotomic extension, with a
specified conjugation and real embedding. State admission requires
Hermiticity, positivity and trace one. Effects satisfy the operator
inequalities between zero and the identity. Some spectral operations
require scalar extensions or supplied witnesses; an exact implementation
can refuse an unsupported input.

One abstract observer model,
\[\protect\adjustbox{max width=\linewidth}{\(\displaystyle \mathcal O=(K,\mathcal F,\mathcal R,\delta,d_0,d_1,J,\mathcal M,\tau),\)}\]
names the carrier, local spaces, comparisons, defects, repairs, closure
test, gauge, compatible transformations and decision rule. For a quantum
experiment it also needs a preparation domain,
system\allowbreak{}/\allowbreak{}apparatus\allowbreak{}/\allowbreak{}environment
split, interaction, instrument, accessible effect menu, readout and
calibration conditions. A physical apparatus realizes part of that
model; a person or conscious experience is not a required mathematical
component.

One declared obstruction observer uses Hermitian edge defects and
\[\protect\adjustbox{max width=\linewidth}{\(\displaystyle J_Q(h)=\frac12\sum_e\|h_e\|_1,
 \qquad
 \Phi_Q([\delta])=\inf_\alpha J_Q(\delta-d_0\alpha).\)}\] For a
finite-dimensional real space of Hermitian edge defects and a linear
repair image, this is the quotient norm induced by \(J_Q\); finite
dimension makes that image closed. It does not identify every formal
repair with a realizable quantum operation. Nor is every algebraic
projector a positive or completely positive channel. Positivity,
physical availability and question preservation must be checked
separately.

Observer relativity therefore means that the experiment and evaluation
are specified. It does not mean that an arbitrary decision label changes
the physical state, or that different descriptions license contradictory
reports about the same compatible experiment.

\subsubsection{13.3 Interaction, readout and conditional
updating}\label{interaction-readout-and-conditional-updating}

On a supplied finite quantum carrier, an instrument has completely
positive maps indexed by record labels. In Kraus form,
\[\protect\adjustbox{max width=\linewidth}{\(\displaystyle \mathcal I_i(\rho)=\sum_a K_{ia}\rho K_{ia}^{\dagger},
 \qquad
 \sum_{i,a}K_{ia}^{\dagger}K_{ia}=I.\)}\] The sum of these maps
preserves trace. Given the positive trace pairing,
\[\protect\adjustbox{max width=\linewidth}{\(\displaystyle E_i=\sum_aK_{ia}^{\dagger}K_{ia},\qquad
 p_i=\operatorname{tr}(E_i\rho)
     =\operatorname{tr}\mathcal I_i(\rho),\qquad
 \rho_i=\mathcal I_i(\rho)/p_i\quad(p_i>0).\)}\] The unnormalized maps
are linear; normalization is conditional on a nonzero-probability
record. A zero-probability branch has no normalized conditional state.
The non-selective state is the sum of the unnormalized branches. These
identities follow from the declared instrument and pairing; they are not
a reconstruction of either from inscription histories.

For an ideal projective instrument, the Kraus operators are the
orthogonal projectors themselves. This gives the Lüders conditional
update and an idempotent non-selective dephasing channel. General
instruments need not be projective, repeatable or idempotent. An effect
specifies an outcome's probability functional, but does not uniquely
specify its disturbance map. The measuring interaction and record
grouping are therefore necessary data.

The exact finite implementation checks these identities and a dilation
for admitted Gaussian-rational inputs. It does not generate one physical
event. Conditional updating answers what follows \textbf{if this record
occurred}; it cannot be used as its own explanation of occurrence.

\subsubsection{13.4 An exact record with recoverable
coherence}\label{an-exact-record-with-recoverable-coherence}

Take two qubits, a system and an apparatus, with
\[\protect\adjustbox{max width=\linewidth}{\(\displaystyle |\psi\rangle=\tfrac35|0\rangle+\tfrac45|1\rangle,
 \qquad |A\rangle=|0\rangle.\)}\] A controlled-NOT from system to
apparatus gives
\[\protect\adjustbox{max width=\linewidth}{\(\displaystyle |\Psi\rangle=\tfrac35|00\rangle+\tfrac45|11\rangle,
 \qquad \rho_{SA}=|\Psi\rangle\langle\Psi|.\)}\] Compare this with the
dephased joint state
\[\protect\adjustbox{max width=\linewidth}{\(\displaystyle \sigma_{SA}=\tfrac9{25}|00\rangle\langle00|
             +\tfrac{16}{25}|11\rangle\langle11|.\)}\] Both states give
the same system reduction, apparatus reduction and pointer distribution:
the diagonal weights are nine twenty-fifths and sixteen twenty-fifths.
Repeated sharp pointer reads are consistent within the ideal instrument.
This is an available classical record distribution under the declared
pointer readout. It is not a proof that the joint state is the mixture
or that one pointer label was selected.

Their difference has only two nonzero entries, each twelve
twenty-fifths, between the basis vectors labelled 00 and 11. Its
eigenvalues are
\[\protect\adjustbox{max width=\linewidth}{\(\displaystyle \tfrac{12}{25},\quad-\tfrac{12}{25},\quad0,\quad0,
 \qquad
 D_{\rm tr}(\rho_{SA},\sigma_{SA})=\tfrac{12}{25}.\)}\] The bounded
observable consisting of Pauli X on each qubit distinguishes them: its
expectation is twenty-four twenty-fifths for the coherent state and zero
for the mixture. Alternatively, applying the same controlled-NOT again
recovers the initial coherent system state from the first preparation;
the mixture yields a diagonal system state with the apparatus reset to
zero. The system's Pauli-X readout then distinguishes them by the same
amount.

Every matrix entry in this example is rational. Its conclusion depends
on the admitted joint controls. If the apparatus or environment cannot
be controlled coherently, this erasure experiment may be unavailable.
Extending the model with an orthogonal environmental record can remove
the coherence from the system\allowbreak{}/\allowbreak{}apparatus
reduction while retaining it globally. Losing access to that distinction
explains the adequacy of a reduced description for a restricted task; it
is not a derivation of objective collapse.

\subsubsection{13.5 When a classical report is
sufficient}\label{when-a-classical-report-is-sufficient}

Let the report be a map from the admissible histories or states to a
record space, and let a protected question have answer map q. Exact
sufficiency is
\[\protect\adjustbox{max width=\linewidth}{\(\displaystyle q=\widehat q\circ p
 \quad\Longleftrightarrow\quad
 p(x)=p(y)\ \Longrightarrow\ q(x)=q(y).\)}\] The factor is unique on the
report image. Equality of the two maps' fibers would impose a stronger
faithfulness condition; factorization only requires that the report
never identifies different protected answers. For a future experiment,
the answer includes its probabilities or conditional reports, so present
agreement alone may be insufficient.

For a supplied nonempty effect menu, denoted separately from the
compatible transformation class by \(\mathcal E_{\rm adm}\), define
accessible distinguishability by
\[\protect\adjustbox{max width=\linewidth}{\(\displaystyle D_{\mathcal E_{\rm adm}}(\rho,\sigma)
   =\sup_{E\in\mathcal E_{\rm adm}}
      |\operatorname{tr}(E(\rho-\sigma))|.\)}\] It is bounded by trace
distance. The full complex effect domain implies equality; a restricted
menu also attains equality if it contains an optimal distinguishing
effect for the given pair. An arbitrary finite menu need not. A
classical description is adequate at a declared tolerance only when it
preserves the protected answers under the admitted preparations, future
controls, readouts and time horizon. If a later control converts a
discarded phase into a readable population, the original population
report was not dynamically sufficient. An algebraic completion may
require observables the apparatus cannot measure jointly or cannot
access at all.

A sharp pointer record algebra generated by orthogonal projectors is
commutative. The complete quantum algebra containing it need not be;
even the block-diagonal fixed-point algebra of dephasing can retain
noncommuting degrees of freedom within a higher-rank block. Thus a
classical record can be a subalgebra of a quantum carrier, rather than a
moment at which the entire carrier becomes classical. Its preferred
basis comes from a specified coupling and stability criterion. A
degenerate coupling does not choose a unique basis without further
apparatus information.

The existing finite sufficient-record derivation
(\texttt{r\allowbreak{}e\allowbreak{}s\allowbreak{}e\allowbreak{}a\allowbreak{}r\allowbreak{}c\allowbreak{}h\allowbreak{}/\allowbreak{}q\allowbreak{}u\allowbreak{}a\allowbreak{}n\allowbreak{}t\allowbreak{}u\allowbreak{}m\allowbreak{}\_\allowbreak{}t\allowbreak{}o\allowbreak{}\_\allowbreak{}c\allowbreak{}l\allowbreak{}a\allowbreak{}s\allowbreak{}s\allowbreak{}i\allowbreak{}c\allowbreak{}a\allowbreak{}l\allowbreak{}\_\allowbreak{}d\allowbreak{}e\allowbreak{}s\allowbreak{}c\allowbreak{}e\allowbreak{}n\allowbreak{}t\allowbreak{}/\allowbreak{}R\allowbreak{}E\allowbreak{}A\allowbreak{}D\allowbreak{}M\allowbreak{}E\allowbreak{}.\allowbreak{}m\allowbreak{}d})
establishes static factorization and invariant-row-space completion on
its declared linear carriers. It also distinguishes autonomous
expectation coordinates from jointly sharp measurements and from a
realized event. Its finite certificates and conditional analytic bounds
retain those premises.

\subsubsection{13.6 Outcomes, observer agreement and the remaining
boundary}\label{outcomes-observer-agreement-and-the-remaining-boundary}

There is a useful elementary restriction on a density-only deterministic
outcome rule. Suppose an affine map sends every density state to a
vertex of a classical probability simplex and sends two pointer
eigenstates to distinct vertices. On their equal mixture, affinity
requires the midpoint of those vertices, which is not a vertex. No such
map satisfies all three requirements. This argument assumes the density
state is the complete input and classical mixing is represented
affinely. It does not exclude an additional configuration, a stochastic
event law, a relative-branch account or another explicitly enlarged
carrier.

Different interpretations provide different further premises. A definite
configuration account supplies variables and dynamics beyond the density
description; an objective-collapse account modifies dynamics; a
relative-state account changes how outcomes are associated with
observers and branches; an operational account can take outcome
occurrence as a primitive. These are distinct positions, not alternative
names for the same derived theorem. None follows solely by assigning a
software label or discarding an environment.

Records from multiple observers require compatible instruments, a
declared comparison protocol and retained history. Separate outcome
tables do not automatically admit a common nonnegative joint assignment
across incompatible contexts. Conversely, contextuality does not prevent
agreement on compatible pointer records. An observer who performs a
coherent erasure has changed the experiment relative to one who
preserves and rereads the mark; a comparison must include that change
rather than merge both into an undeclared history.

Our current answer therefore explains, under supplied quantum premises,
how interactions furnish record probabilities, why particular reduced
descriptions can be adequate and stable, and how conditional updates
compose. It leaves two foundational burdens explicit: deriving the
operational carrier from native histories, and accounting for actual
outcome occurrence under a stated ontology or event law. Both questions
remain open.

\subsubsection{13.7 Testing two foundational carrier
types}\label{testing-two-foundational-carrier-types}

Begin with the hypothesis that a quantum carrier and a classical carrier
are distinct foundational types. The claim becomes testable only after
declaring their state spaces, readouts, allowed transfers and
composition. Calling the types distinct does not specify a prediction.
Conversely, embedding a classical register into a quantum matrix algebra
does not prove that only quantum things exist. We separate the following
versions of the hypothesis.

\begin{longtable}[]{@{}
  >{\raggedright\arraybackslash}p{(\linewidth - 4\tabcolsep) * \real{0.3333}}
  >{\raggedright\arraybackslash}p{(\linewidth - 4\tabcolsep) * \real{0.3333}}
  >{\raggedright\arraybackslash}p{(\linewidth - 4\tabcolsep) * \real{0.3333}}@{}}
\toprule\noalign{}
\begin{minipage}[b]{\linewidth}\raggedright
Hypothesis
\end{minipage} & \begin{minipage}[b]{\linewidth}\raggedright
Finite operational content
\end{minipage} & \begin{minipage}[b]{\linewidth}\raggedright
Result here
\end{minipage} \\
\midrule\noalign{}
\endhead
\bottomrule\noalign{}
\endlastfoot
Two types coexist & Quantum systems and classical registers have their
own admitted states and transfers. & Consistent hybrid models exist; the
bare claim is not refuted. \\
A population record is the complete future carrier & Presently equal
populations imply equal answers after every admitted control. & Refuted
on the declared coherent-control protocol below. \\
Classical storage faithfully replaces quantum storage & Positive,
normalized encoding and retrieval preserve arbitrary unknown qubit
inputs. & Refuted by the sharp four-state bound below. \\
Additional classical variables explain quantum predictions & An enlarged
classical description may depend on preparation, context, hidden
variables or nonlocal dynamics. & Not excluded by the storage bound; its
own premises and predictions must be tested. \\
The full quantum carrier is all that exists & No additional ontology is
needed or exists. & Operational sufficiency does not decide this
ontological assertion. \\
\end{longtable}

\textbf{A consistent two-type countermodel.} Let
\[\protect\adjustbox{max width=\linewidth}{\(\displaystyle \mathcal A=\bigoplus_{\lambda=1}^{m}M_2(\mathbb C).\)}\]
A state is a tuple \((\rho_\lambda)_\lambda\), with each block positive
and
\[\protect\adjustbox{max width=\linewidth}{\(\displaystyle \sum_\lambda\operatorname{tr}\rho_\lambda=1.\)}\]
Effects are tuples \(0\le E_\lambda\le I\); their probabilities are
\[\protect\adjustbox{max width=\linewidth}{\(\displaystyle \sum_\lambda\operatorname{tr}(E_\lambda\rho_\lambda).\)}\]
The center carries a classical label, while each block retains
noncommuting qubit observables. Stochastic label changes and
within-block quantum channels furnish consistent finite operations. A
normalized POVM and supplied density preparations give the block
transfer
\[\protect\adjustbox{max width=\linewidth}{\(\displaystyle \rho_\lambda=\operatorname{tr}(F_\lambda\rho)\sigma_\lambda.\)}\]
Summing the blocks discards the label. The model can be presented over
\(\mathbb Q(i)\) on its exact finite fixtures.

This is a genuine restriction of the full \(M_{2m}\) quantum model if
coherences and controls between label blocks are prohibited. Its matrix
embedding does not remove that restriction or decide its ontology. It
also supplies a distribution on labels, not a derivation of one
occurring event. Thus the existence of two operational types is
mathematically consistent. A claim that a particular apparatus must
inhabit such a restricted type needs evidence for the prohibited
coherences and controls.

\textbf{Present equality is not future equality.} In the declared qubit
basis let
\[\protect\adjustbox{max width=\linewidth}{\(\displaystyle |\psi_\pm\rangle=\tfrac35|0\rangle\pm\tfrac45|1\rangle,
 \qquad
 U=\begin{pmatrix}3/5&4/5\\-4/5&3/5\end{pmatrix}.\)}\] Both states have
population report \((9/25,16/25)\), but after the same available unitary
\(U\), the probability of the first basis outcome is respectively \(1\)
and \(49/625\). A model that identifies these states as the same
complete state cannot predict both answers using only that record.
Dephasing them to the same diagonal density operator makes their future
answers equal; it has changed the state, rather than faithfully retained
the discarded phase. This establishes that the population report is
insufficient for that future question. It does not rule out classical
storage of additional variables, nor assume that the displayed control
is available in every device.

\textbf{A sharp classical-storage obstruction.} Consider an
unconditional, trace-preserving transfer of an unknown input qubit
through a finite classical alphabet and back to a qubit. The
intermediate outcome may be random; the protocol must succeed without
postselection. Use the full positive qubit
state\allowbreak{}/\allowbreak{}effect model and affine mixing. The
encoder sees only the input state; it has no preparation label, quantum
side memory or retained coherent environment. The decoder returns a
normalized positive qubit state for each letter. No postselection or
conditional success loophole is included. The test basis is disclosed
only after the returned qubit is fixed; requesting a basis-dependent
classical answer is a different protocol. These premises give the
measure-and-prepare form
\[\protect\adjustbox{max width=\linewidth}{\(\displaystyle \Phi(\rho)=\sum_\lambda
     \operatorname{tr}(F_\lambda\rho)\sigma_\lambda,
 \quad F_\lambda\ge0,\quad\sum_\lambda F_\lambda=I,
 \quad\sigma_\lambda\ge0,\quad
     \operatorname{tr}\sigma_\lambda=1.\)}\] A positive affine encoder
on the full density-state domain has POVM response operators; an affine
decoder from a probability simplex is specified by its states at the
vertices. Thus the displayed form follows from the stated operational
premises. It is also the standard classical-transfer form of an
entanglement-breaking channel {[}11{]}. The finite Choi matrix, with
input factor first, is
\[\protect\adjustbox{max width=\linewidth}{\(\displaystyle J_\Phi=\sum_\lambda F_\lambda^T\otimes\sigma_\lambda.\)}\]
It is positive, and tracing its output factor gives \(I\); complete
positivity and trace preservation therefore follow without choosing
Kraus square roots outside the coefficient field.

Let the four equiprobable inputs be
\[\protect\adjustbox{max width=\linewidth}{\(\displaystyle P_{Z\pm}=\tfrac12(I\pm Z),\qquad
 P_{X\pm}=\tfrac12(I\pm X),\)}\] and define their average survival
probability
\[\protect\adjustbox{max width=\linewidth}{\(\displaystyle S(\Phi)=\frac14\sum_{a\in\{X,Z\},\,s\in\{+,-\}}
      \operatorname{tr}\bigl(P_{as}\Phi(P_{as})\bigr).\)}\] All four
projectors have rational entries. The exact theorem is
\[\protect\adjustbox{max width=\linewidth}{\(\displaystyle \boxed{S(\Phi)\le\frac34.}\)}\]
It holds for every finite number of classical letters under the premises
above, and the constant is attained.

For a proof, use the real trace coordinates
\[\protect\adjustbox{max width=\linewidth}{\(\displaystyle w_\lambda=\operatorname{tr}F_\lambda,\qquad
 f_{\lambda a}=\operatorname{tr}(F_\lambda a),\qquad
 s_{\lambda a}=\operatorname{tr}(\sigma_\lambda a).\)}\] The two-by-two
determinant and trace conditions for positivity imply
\[\protect\adjustbox{max width=\linewidth}{\(\displaystyle f_{\lambda X}^{\,2}+f_{\lambda Z}^{\,2}\le w_\lambda^2,
 \qquad
 s_{\lambda X}^{\,2}+s_{\lambda Z}^{\,2}\le1.\)}\] The omitted \(Y\)
coordinates can only strengthen these inequalities. Cauchy--Schwarz and
\(w_\lambda\ge0\) give the nonnegative margin
\[\protect\adjustbox{max width=\linewidth}{\(\displaystyle d_\lambda=w_\lambda
       -f_{\lambda X}s_{\lambda X}
       -f_{\lambda Z}s_{\lambda Z}\ge0.\)}\] Expanding the four traces
and using \(\sum_\lambda w_\lambda=2\) yields
\[\protect\adjustbox{max width=\linewidth}{\(\displaystyle S(\Phi)=\frac12+\frac{1}{8}\sum_\lambda
      (f_{\lambda X}s_{\lambda X}+f_{\lambda Z}s_{\lambda Z}).\)}\]
\[\protect\adjustbox{max width=\linewidth}{\(\displaystyle \frac34-S(\Phi)=\frac18\sum_\lambda d_\lambda.\)}\]
Measuring \(Z\) and preparing its obtained eigenstate gives survival
probabilities \((1,1,1/2,1/2)\) in the order \(Z+,Z-,X+,X-\), attaining
\(S=3/4\). Identity quantum storage gives \(S=1\). A faithful positive
classical replacement preserving even these four preparations and their
retrieval tests therefore cannot exist under these premises. More
classical letters, shared randomness independent of the input, or an
additional letter cannot improve the bound: each normalized strategy has
the same inequality, and convex mixtures retain it.

For a declared error criterion, if all four survival probabilities are
at least \(1-\varepsilon\), the bound requires \(\varepsilon\ge1/4\).
More generally, if each retrieved state is within trace distance
\(\varepsilon\) of its input, its own projector's survival probability
is at least \(1-\varepsilon\), giving the same necessary bound. Randomly
choosing between \(X\) and \(Z\)
measurement\allowbreak{}/\allowbreak{}preparation with equal weights
attains survival \(3/4\) and trace distance \(1/4\) on each of the four
states: the retrieved state is \(\tfrac34P+\tfrac14(I-P)\), whose
difference from \(P\) has eigenvalues \(\pm1/4\). Both four-state error
bounds are sharp. This is a bound on an unconditional classical-storage
protocol, not a noise estimate for an arbitrary apparatus or a statement
about selectively successful trials.

\textbf{A static classical control and the boundary of the conclusion.}
Four classical letters with predetermined signs \((z,x)\in\{+,-\}^2\)
reproduce all the four preparations' static \(X/Z\) readouts: for a
\(Z\) eigenstate, put equal probability on the two matching \(z\)
letters, and similarly for an \(X\) eigenstate. This also preserves the
equal-mixture identity
\[\protect\adjustbox{max width=\linewidth}{\(\displaystyle \tfrac12(P_{Z+}+P_{Z-})=\tfrac12(P_{X+}+P_{X-})=I/2.\)}\]
These finite tables do not refute a classical explanation.

What they cannot supply is a positive
encoding\allowbreak{}/\allowbreak{}retrieval on the full qubit carrier.
A letter with definite \(x=z=+1\) would require a retrieved qubit with
both expectations equal to one: the unit-probability readouts force that
requirement on every letter with nonzero weight in both preparations.
The candidate \((I+X+Z)/2\) has determinant \(-1/4\); allowing an extra
\(Y\) coefficient makes this determinant still smaller. Similarly,
extending that letter's response probability affinely to all qubits
requires \(F_{++}=(I+X+Z+bY)/4\), whose determinant is \(-(1+b^2)/16\).
Positivity fails. Restricting the admitted states and probes to the
four-state convex hull, allowing contextual responses, or supplying a
preparation label changes these premises and can evade the conclusion. A
classical computer storing known matrix entries is also a description of
inputs, not a positive physical encoding of an unknown qubit.

The result establishes a concrete asymmetry: classical records can live
inside a supplied quantum model, and may be sufficient for selected
questions; an unknown quantum state cannot in general be replaced and
restored through a solely classical positive carrier. Classical-looking
behavior can reflect restricted access to coherence, actual loss of
coherence, or a restricted hybrid model. Those possibilities must be
distinguished by declared future controls and protected questions.
Neither this obstruction nor coherent erasure refutes every two-type
ontology, establishes that quantum ontology is exhaustive, or supplies
the missing event-selection law.

\subsubsection{13.8 Three separate carrier
distinctions}\label{three-separate-carrier-distinctions}

The two-type hypothesis becomes sharper when three different
obstructions are kept separate. A report can lose a required future
answer; compatible local probability tables can fail to have a positive
global extension; and a declared classical center can conflict with
purification of every state. None of these statements alone establishes
that the quantum carrier is all that exists.

\textbf{Test 1: future-question sufficiency.} Supply a finite
preparation family, a present POVM report, a CPTP control and a future
POVM. The report is the whole ensemble probability vector, not one
sampled outcome. Ask whether the probability of outcome zero of the
future POVM factors through that report. Two preparations with the same
report and different answers disprove such a decoder; otherwise a finite
decoder is returned on the supplied family. For
\[\protect\adjustbox{max width=\linewidth}{\(\displaystyle |\psi_\pm\rangle=(3|0\rangle\pm4|1\rangle)/5,\qquad
 U=\frac15\begin{pmatrix}3&4\\-4&3\end{pmatrix},\)}\] Z populations
coincide at \((9/25,16/25)\), but the probabilities of zero after \(U\)
are \(1\) and \(49/625\). Identity control passes the same sufficiency
test. Reporting the pulled-back future effect also restores sufficiency
on this family. This is the finite protected-question criterion and the
phase-blind descent obstruction on an exact fixture. A pass on these two
states does not prove sufficiency for every preparation or every future
control.

\textbf{Test 2: positive Bell gluing.} Supply the four binary pair
distributions for settings \(00,01,10,11\). Ask for nonnegative weights
on the sixteen assignments \((a_0,a_1,b_0,b_1)\) reproducing all four
marginals. Exact convex membership returns either those weights or a
separating functional verified on all sixteen deterministic vertices.
Signed solutions to linear equations are insufficient. For
\(|\Phi^+\rangle\), take
\[\protect\adjustbox{max width=\linewidth}{\(\displaystyle A_0=Z,\quad A_1=X,\quad B_0=(4Z+3X)/5,\quad B_1=(4Z-3X)/5.\)}\]
The correlation vector is \((4/5,4/5,3/5,-3/5)\), hence
\[\protect\adjustbox{max width=\linewidth}{\(\displaystyle S=E_{00}+E_{01}+E_{10}-E_{11}=14/5>2\)}\].
Every context probability is strictly positive: their numerators over 20
are respectively
\[\protect\adjustbox{max width=\linewidth}{\(\displaystyle (9,1,1,9),(9,1,1,9),(8,2,2,8),(2,8,8,2)\)}\],
in \((++,+-,-+,--)\) order. Thus a support-only compatibility test
admits all sixteen assignments and misses this obstruction. Another
regression has \(S=0\) for this orientation but violates a different
CHSH facet; the implementation tests the complete convex hull. Uniform
independent outcomes provide the positive extension control. The
numerators displayed here are the ones encoded in the Lean fixture; the
Rust violation door V2 in source-notes supplement §11.33 uses the
companion orientation \(B_1=(3Z-4X)/5\), with correlation vector
\((4/5,3/5,3/5,-4/5)\) and the same \(S=14/5\).

Bell locality and setting independence explain why one would require
such a joint table in a physical model; they are additional premises,
not outputs of the solver. In detail, suppose a common,
setting-independent probability measure \(\mu\) on hidden states
\(\lambda\) and local response distributions \(A(a\mid x,\lambda)\),
\(B(b\mid y,\lambda)\) give the four pair tables. Then
\[\protect\adjustbox{max width=\linewidth}{\(\displaystyle w(a_0,a_1,b_0,b_1)=\int
 A(a_0\mid0,\lambda)A(a_1\mid1,\lambda)
 B(b_0\mid0,\lambda)B(b_1\mid1,\lambda)\,d\mu(\lambda)\)}\] is a
nonnegative normalized table on the sixteen assignments and has those
pair marginals: sum out the unused response factors, whose sums are one.
This reduction works for finite, countable or continuous hidden-state
spaces, provided the stated probabilities are measurable. Increasing
their cardinality does not evade the obstruction. This finite result
excludes that conjunction for the supplied statistics. It does not
exclude nonlocal hidden-variable models or establish a unique ontology.
The obstruction is failure of positivity-constrained global gluing,
which must not be conflated with signed linear exactness or with a
support-only contextuality obstruction. The source manuscript on finite
Bell bounds (source-notes supplement §11.33) already identifies this
joint-table premise; this test adds a shared executable experiment
{[}26{]} and a Lean marginal-to-CHSH proof.

\textbf{Test 3: central sectors and purification.} Declare the full
positive normalized states of \(A=\bigoplus_i M_{n_i}(\mathbb C)\),
standard tensor composition with another such algebra, and the available
ancilla block sizes. A state has normalized central weights \(p_i\) and
block states \(\rho_i\). A pure composite state occupies one joint
central block, so its marginal has one central weight equal to one. A
state with two positive central weights therefore has no pure extension
within this model, for any ancilla of this form. More quantitatively,
every extension \(\omega\) of that state and every marginal \(\sigma\)
of a pure composite state obey
\[\protect\adjustbox{max width=\linewidth}{\(\displaystyle \operatorname{Tr}(\omega^2)\le\sum_i p_i^2,\qquad
 \tfrac12\|\rho-\sigma\|_1\ge1-\max_i p_i.\)}\] For the first
inequality, sum the squared weights of normalized joint blocks, use
purity at most one within each block, and then use
\[\protect\adjustbox{max width=\linewidth}{\(\displaystyle \sum_j p_{ij}^2\le(\sum_j p_{ij})^2=p_i^2\)}\].
For the second, the central measurement is a trace-distance contraction
and distinguishes \(p\) from a point mass by \(1-p_k\). With a
sufficiently large quantum ancilla both bounds are sharp: purify each
\(\rho_i\) inside a chosen ancilla block for the first; for the second
purify \(\rho_k\) in a block of maximal central weight. The half-half
classical mixture gives both bounds \(1/2\). The exact finite
calculation separately computes tighter purity, rank and
marginal-distance constraints for the listed ancilla sizes.

A single sector has no central obstruction, but an undersized ancilla
can still prevent purification. The control with spectrum
\((9/25,16/25)\) needs an ancilla block of size at least two. Requiring
purification for every state rules out nontrivial centers only under
these state-space, composition and purity premises. It neither rules out
a consistent hybrid model that declines that requirement nor derives
purification as a law of nature. This is the conditional
center-elimination step in the companion obstruction analysis; real
quantum theory's purification and the associativity of quaternion
multiplication are unaffected.

\paragraph{13.8.1 Relation to finite report
quotients}\label{relation-to-finite-report-quotients}

The three tests address distinct interfaces. Phase-sensitive readouts
ask whether a report preserves answers to declared future experiments.
Positive Bell extension asks whether compatible pair distributions are
marginals of one nonnegative joint law. The central-sector bounds ask
what extensions are possible in a specified direct-sum tensor model.
None of these questions selects the underlying quantum state space.

A finite protected-question quotient identifies two histories only when
every admitted continuation gives the same answer and remains available.
That is a minimal report relative to one finite experiment, not a
derivation of positive operator algebras, complex scalars, or tensor
composition. Similarly, true marginals and products of one-site
marginals can each satisfy their consistency conditions even when their
difference is nonzero. Report locality, local tomography, and Bell local
causality therefore require different hypotheses. The detailed
comparison and local source addresses are in the source-notes supplement
and companion manuscript {[}12{]}.

\paragraph{13.8.2 From primitive difference to Bell: the endpoint and
the open
step}\label{from-primitive-difference-to-bell-the-endpoint-and-the-open-step}

Two questions must be separated. First, the finite algebra of a
Bell\allowbreak{}/\allowbreak{}CHSH model closes once its state,
observables, composition and probability rule are supplied. Second,
primitive difference alone does not force those quantum inputs. Positive
\(L^1\) weights are additional carrier data, not a theorem of a bare
distinction relation. Once four binary response variables
\[\protect\adjustbox{max width=\linewidth}{\(\displaystyle A_0,A_1,B_0,B_1\in\{\mathord{-}1,+1\}\)}\]
share one positive normalized joint law, every deterministic assignment
has

\[\protect\adjustbox{max width=\linewidth}{\(\displaystyle S=A_0(B_0+B_1)+A_1(B_0-B_1)\in\{-2,2\},\)}\]

and convexity gives \(|S|\leq2\). This conditional joint-extension
theorem is proved in source-notes supplement §11.33; it is not a
derivation of Bell's theorem from \(L^1\) alone. In the CHSH scenario,
Fine's theorem identifies existence of the relevant common joint
distribution with satisfaction of the complete
Bell\allowbreak{}/\allowbreak{}CHSH family under the compatibility
assumptions {[}16{]}. Bell's physical conclusion additionally uses
locality and an identification of preparations, settings and outcomes
{[}14,15{]}.

The quantum side is now equally explicit. On the declared bipartite
carrier, let

\[\protect\adjustbox{max width=\linewidth}{\(\displaystyle \begin{aligned}
A_0&=Z,& A_1&=X,\\
B_0&=(Z+X)/\sqrt2,& B_1&=(Z-X)/\sqrt2,
\end{aligned}\)}\]

and use the entangled state
\[\protect\adjustbox{max width=\linewidth}{\(\displaystyle |\Phi^+\rangle=(|00\rangle+|11\rangle)/\sqrt2\)}\]
with the trace\allowbreak{}/\allowbreak{}Born pairing. The exact
implementation constructs the CHSH operator and verifies the two
sum-of-squares inequalities

\[\protect\adjustbox{max width=\linewidth}{\(\displaystyle -2\sqrt2 I\leq B_{\rm CHSH}\leq2\sqrt2 I\, .\)}\]

The same witness evaluates to \(S=2\sqrt2\) {[}17{]}. Every matrix entry
in this witness is fixed by complex conjugation: the construction lives
in the real field \(\mathbb Q(\sqrt2)\). It is also static. Therefore a
specifically complex \(U(1)\) phase fiber and continuous unitary time
evolution are not necessary for a Tsirelson-saturating CHSH witness.

\begin{longtable}[]{@{}
  >{\raggedright\arraybackslash}p{(\linewidth - 4\tabcolsep) * \real{0.3333}}
  >{\raggedright\arraybackslash}p{(\linewidth - 4\tabcolsep) * \real{0.3333}}
  >{\raggedright\arraybackslash}p{(\linewidth - 4\tabcolsep) * \real{0.3333}}@{}}
\toprule\noalign{}
\begin{minipage}[b]{\linewidth}\raggedright
Boundary
\end{minipage} & \begin{minipage}[b]{\linewidth}\raggedright
What is now closed
\end{minipage} & \begin{minipage}[b]{\linewidth}\raggedright
What remains supplied or open
\end{minipage} \\
\midrule\noalign{}
\endhead
\bottomrule\noalign{}
\endlastfoot
primitive difference to positive weights & finite
difference\allowbreak{}/\allowbreak{}cochain and quotient machinery &
positivity and normalized preparation weights are admitted data \\
classical Bell face & one positive global joint law implies
\(\lvert S\rvert\leq2\) exactly & physical locality, setting
independence and experiment identification \\
strict quantum violation & rational density matrix and observables give
\(S=14/5>2\) and no positive joint table & tensor composition, entangled
preparation, measurements and Born\allowbreak{}/\allowbreak{}trace
readout \\
quantum endpoint & exact SOS upper bound and exact witness give
\(S=2\sqrt2\) over \(\mathbb Q(\sqrt2)\) & derivation or operational
selection of the quantum inputs from native histories \\
\end{longtable}

This also corrects the proposed ``five-import boundary.'' Its five
imports (the phase fiber, the quadratic readout, context parity, the
frame function and the clock tower, each mapped to finite data in a
research note
\texttt{r\allowbreak{}e\allowbreak{}s\allowbreak{}e\allowbreak{}a\allowbreak{}r\allowbreak{}c\allowbreak{}h\allowbreak{}/\allowbreak{}q\allowbreak{}m\allowbreak{}\_\allowbreak{}2\allowbreak{}0\allowbreak{}7\allowbreak{}\_\allowbreak{}*}
to
\texttt{r\allowbreak{}e\allowbreak{}s\allowbreak{}e\allowbreak{}a\allowbreak{}r\allowbreak{}c\allowbreak{}h\allowbreak{}/\allowbreak{}q\allowbreak{}m\allowbreak{}\_\allowbreak{}2\allowbreak{}1\allowbreak{}1\allowbreak{}\_\allowbreak{}*})
remain useful as a catalogue of structures used by \textbf{full complex
quantum mechanics}, but they are neither proved independent nor a
minimal basis for Bell.
Quadratic\allowbreak{}/\allowbreak{}inner-product geometry and the Born
pairing overlap rather than forming an established independence theorem.
Same-party noncommutativity is load-bearing for the maximal CHSH
construction, while cross-party observables commute. Bipartite
composition and entangled preparation are load-bearing but absent from
the five-item cascade. Complex phase and continuous dynamics belong to
the larger theory, not to the static Bell minimum.

\textbf{Verdict.} The finite algebraic road ends at Bell: conditional on
the declared quantum interface, this repository now contains the
classical bound, an exact violation, the Tsirelson upper certificate and
an attaining witness. The foundational road from primitive difference
does not end there. Its remaining step is the first derivation target in
Source Notes §6 (retained subsection §13.9): derive or operationally
justify the positive composite algebra, state, effects and pairing
without silently importing the quantum model that the Bell calculation
uses.

\paragraph{13.8.3 Čech contextuality and the Regge
firewall}\label{ux10dech-contextuality-and-the-regge-firewall}

The sheaf-theoretic language sharpens the local-to-global part of this
story, but only when its objects and degrees are kept distinct. A
measurement scenario consists of a measurement set and a cover by
compatible contexts. An empirical model is a family of probability
distributions on those contexts whose marginals agree on overlaps. In
Bell scenarios this compatibility is the no-signalling condition; in
general contextuality scenarios it is often called no-disturbance. This
is the compatible-family, or Čech degree-zero, condition for the
declared distribution presheaf {[}18{]}. It is not yet a
first-cohomology obstruction.

Abramsky and Brandenburger identify a global section of the distribution
presheaf with a noncontextual hidden-variable realization for a general
measurement cover {[}18{]}. Fine's theorem is the more specific CHSH
result used in §13.8.2 {[}16{]}. The cohomological refinement instead
chooses a local section in an abelian presheaf derived from the support
and constructs an obstruction class. A nonzero class is sufficient for
contextuality, but the test is not complete: contextual models with a
vanishing chosen cohomological obstruction exist {[}19{]}. Three
distinctions follow.

\begin{enumerate}
\def\labelenumi{\arabic{enumi}.}
\item
  Vanishing of the relevant obstruction class does not say that the
  whole cohomology group is zero.
\item
  Compatible local distributions do not automatically constitute a
  nonzero Čech \(H^1\) class; the coefficient presheaf, chosen local
  section and coboundary maps must be declared.
\item
  The scalar CHSH excess

  \[\protect\adjustbox{max width=\linewidth}{\(\displaystyle \Delta_{\rm CHSH}=\max(0,|S|-2)\)}\]

  is a useful violation margin, not a magnitude or norm of a Čech class.
  Thus the rational witness in §13.8.2 has \(\Delta_{\rm CHSH}=4/5\),
  and a PR box would have \(\Delta_{\rm CHSH}=2\), but these arithmetic
  facts do not define a cohomology theory. The rational witness excludes
  a positive global joint law; it does not, without a separately
  declared support construction, compute a Čech obstruction. PR boxes
  are among the examples for which the standard support-derived
  obstruction is known to be nonzero {[}19{]}.
\end{enumerate}

The Mermin--Peres scenario does support an exact finite local-to-global
obstruction. Let
\[\protect\adjustbox{max width=\linewidth}{\(\displaystyle d_0:\mathbb F_2^V\to\mathbb F_2^C\)}\]
map nine observable values to the six context parities, let
\(s\in\mathbb F_2^C\) be the declared sign vector, and let \(\ell\) sum
all six context coordinates. Every observable occurs twice, so
\(\ell d_0=0\), while the odd target has \(\ell(s)=1\). Hence
\(s\notin\operatorname{im}d_0\). Every individual context and every
five-context subfamily is satisfiable, but the complete six-context
family is not. This is the exact locally satisfiable but globally
inconsistent parity obstruction verified in source-notes supplement
§11.50 and is the finite combinatorial core of the Peres--Mermin
argument {[}5,6,20{]}. It is safe to call it a nonzero cokernel class.
Calling it an \(H^1\) class additionally requires a declared next
coboundary \(d_1\) with \(d_1s=0\); calling it a class in \(H^2\)
requires a declared two-complex. The bare carrier graph has no
degree-two cochains. Finally, the parity contradiction itself is finite
\(\mathbb F_2\) algebra, but identifying its signs with quantum
measurements uses the supplied commuting two-qubit Pauli operator
realization. It is not a derivation of that quantum realization from
parity alone.

Phase holonomy needs the same firewall. Transition functions for a
specified \(U(1)\) bundle can form a Čech 1-cocycle, but Aharonov--Bohm
and Berry holonomies also depend on a connection and a loop {[}22,23{]}.
Nontrivial connection holonomy can occur on a trivial bundle, so it is
not automatically a nonzero class in a bare \(\check H^1\) of transition
data. Likewise, calling a phase a coboundary requires a declared
coefficient sheaf, cover and coboundary map. The paper therefore retains
phase holonomy as a conditional structural analogy, not as a general
identification with the contextuality obstruction above.

The same discipline blocks the proposed Regge identification. Regge
calculus places curvature on codimension-two hinges, with deficit

\[\protect\adjustbox{max width=\linewidth}{\(\displaystyle \epsilon_h=2\pi-\sum_{\sigma\supset h}\theta_{\sigma h}\)}\]

for the appropriate dihedral angles {[}21{]}. In a two-dimensional
equilateral triangulation the hinge is a vertex: six triangles give zero
deficit, five a positive deficit and seven a negative deficit. That is a
geometric curvature diagnostic. It neither supplies a measurement cover
nor empirical probabilities. Even zero local deficits do not, without
global hypotheses, force trivial topology or holonomy.

Accordingly, no result here identifies a Čech obstruction with
\(|\epsilon_h|\), establishes the bound
\[\protect\adjustbox{max width=\linewidth}{\(\displaystyle S_{\rm CHSH}\leq\min(4,2+|\epsilon_h|)\)}\]
or makes curvature a prerequisite for Bell violation. The existing
experimental API records the angle but returns
\texttt{N\allowbreak{}o\allowbreak{}n\allowbreak{}e} for the Čech
obstruction, Bell bound and contextuality flag, including at zero
deficit. The exact CHSH witness in §13.8.2 is independent of Regge
curvature and reaches \(2\sqrt2\), so the proposed unconditional
implication \(\epsilon=0\Rightarrow |S|\leq2\) cannot be a general
quantum theorem.

\textbf{Narrow claim and next test.} Čech cohomology can witness a
declared local-to-global obstruction, while Regge deficits measure
declared simplicial curvature; their resemblance is presently an
analogy, not a bridge. A genuine bridge must map a specified
triangulation and hinge data to a measurement cover, outcome presheaf
and empirical distributions, prove overlap compatibility, define the
coefficient complex, and then derive a bound that survives both a
classical control and the zero-deficit quantum CHSH counterexample.
Until those data and proofs exist, the proposed bridge remains open.

\subsection{References}\label{references}

Numbers in square brackets refer to this list; it is also shipped as
\texttt{d\allowbreak{}o\allowbreak{}c\allowbreak{}s\allowbreak{}/\allowbreak{}r\allowbreak{}e\allowbreak{}f\allowbreak{}e\allowbreak{}r\allowbreak{}e\allowbreak{}n\allowbreak{}c\allowbreak{}e\allowbreak{}s\allowbreak{}.\allowbreak{}m\allowbreak{}d}.

\begin{enumerate}
\def\labelenumi{\arabic{enumi}.}
\tightlist
\item
  A. M. Gleason, ``Measures on the Closed Subspaces of a Hilbert
  Space,'' Journal of Mathematics and Mechanics 6, 885--893 (1957).
\item
  A. Cabello, J. M. Estebaranz and G. García-Alcaine,
  ``Bell-Kochen-Specker Theorem: A Proof with 18 Vectors,'' Physics
  Letters A 212, 183--187 (1996).
\item
  P. Busch, ``Quantum States and Generalized Observables: A Simple Proof
  of Gleason's Theorem,'' Physical Review Letters 91, 120403 (2003).
\item
  H. Weyl, ``The Theory of Groups and Quantum Mechanics,'' Dover
  Publications (1950).
\item
  N. D. Mermin, ``Hidden Variables and the Two Theorems of John Bell,''
  Reviews of Modern Physics 65, 803--815 (1993).
\item
  A. Peres, ``Incompatible Results of Quantum Measurements,'' Physics
  Letters A 151, 107--108 (1990).
\item
  R. D. Sorkin, ``Quantum Mechanics as Quantum Measure Theory,'' Modern
  Physics Letters A 9, 3119--3127 (1994).
\item
  M. Schlosshauer, ``Decoherence, the Measurement Problem, and
  Interpretations of Quantum Mechanics,'' Reviews of Modern Physics 76,
  1267--1305 (2004).
  \href{https://arxiv.org/abs/quant-ph/0312059}{arXiv:quant-ph\allowbreak{}/\allowbreak{}0312059}.
\item
  T. Maudlin, ``Three Measurement Problems,'' Topoi 14, 7--15 (1995).
  \href{https://doi.org/10.1007/BF00763473}{\nolinkurl{doi:10.1007/BF00763473}}.
\item
  W. H. Zurek, ``Decoherence, Einselection, and the Quantum Origins of
  the Classical,'' Reviews of Modern Physics 75, 715--775 (2003).
  \href{https://doi.org/10.1103/RevModPhys.75.715}{\nolinkurl{doi:10.1103/RevModPhys.75.715}}.
\item
  M. Horodecki, P. W. Shor and M. B. Ruskai, ``General Entanglement
  Breaking Channels'' (2003).
  \href{https://arxiv.org/abs/quant-ph/0302031}{arXiv:quant-ph\allowbreak{}/\allowbreak{}0302031};
  journal version
  \href{https://doi.org/10.1142/S0129055X03001709}{\nolinkurl{doi:10.1142/S0129055X03001709}}.
\item
  J. H. Carroll, \emph{Projection Obstruction for Factorized States},
  author-review version 6.9.16 (2026), §§4.3.2, 7.6, 7.10 and Appendix
  E. Local companion address:
  \texttt{c\allowbreak{}o\allowbreak{}m\allowbreak{}p\allowbreak{}a\allowbreak{}n\allowbreak{}i\allowbreak{}o\allowbreak{}n\allowbreak{}\_\allowbreak{}p\allowbreak{}a\allowbreak{}p\allowbreak{}e\allowbreak{}r\allowbreak{}s\allowbreak{}/\allowbreak{}p\allowbreak{}r\allowbreak{}o\allowbreak{}j\allowbreak{}e\allowbreak{}c\allowbreak{}t\allowbreak{}i\allowbreak{}o\allowbreak{}n\allowbreak{}\_\allowbreak{}o\allowbreak{}b\allowbreak{}s\allowbreak{}t\allowbreak{}r\allowbreak{}u\allowbreak{}c\allowbreak{}t\allowbreak{}i\allowbreak{}o\allowbreak{}n\allowbreak{}.\allowbreak{}m\allowbreak{}d}.
\item
  J. H. Carroll, \emph{Trace-Zero Repair Quotients on a Finite Graph: An
  Exact Qutrit Counterexample and the Correct Scalar-Divergence Dual},
  Version 0.55.0 (2026). Local companion address:
  \texttt{c\allowbreak{}o\allowbreak{}m\allowbreak{}p\allowbreak{}a\allowbreak{}n\allowbreak{}i\allowbreak{}o\allowbreak{}n\allowbreak{}\_\allowbreak{}p\allowbreak{}a\allowbreak{}p\allowbreak{}e\allowbreak{}r\allowbreak{}s\allowbreak{}/\allowbreak{}t\allowbreak{}r\allowbreak{}a\allowbreak{}c\allowbreak{}e\allowbreak{}\_\allowbreak{}z\allowbreak{}e\allowbreak{}r\allowbreak{}o\allowbreak{}\_\allowbreak{}q\allowbreak{}u\allowbreak{}o\allowbreak{}t\allowbreak{}i\allowbreak{}e\allowbreak{}n\allowbreak{}t\allowbreak{}\_\allowbreak{}d\allowbreak{}u\allowbreak{}a\allowbreak{}l\allowbreak{}i\allowbreak{}t\allowbreak{}y\allowbreak{}\_\allowbreak{}m\allowbreak{}o\allowbreak{}n\allowbreak{}o\allowbreak{}g\allowbreak{}r\allowbreak{}a\allowbreak{}p\allowbreak{}h\allowbreak{}.\allowbreak{}m\allowbreak{}d}.
\item
  J. S. Bell, ``On the Einstein Podolsky Rosen Paradox,'' \emph{Physics}
  \textbf{1}, 195--200 (1964).
  \href{https://doi.org/10.1103/PhysicsPhysiqueFizika.1.195}{\nolinkurl{doi:10.1103/PhysicsPhysiqueFizika.1.195}}.
\item
  J. F. Clauser, M. A. Horne, A. Shimony and R. A. Holt, ``Proposed
  Experiment to Test Local Hidden-Variable Theories,'' \emph{Physical
  Review Letters} \textbf{23}, 880--884 (1969).
  \href{https://doi.org/10.1103/PhysRevLett.23.880}{\nolinkurl{doi:10.1103/PhysRevLett.23.880}}.
\item
  A. Fine, ``Hidden Variables, Joint Probability, and the Bell
  Inequalities,'' \emph{Physical Review Letters} \textbf{48}, 291--295
  (1982).
  \href{https://doi.org/10.1103/PhysRevLett.48.291}{\nolinkurl{doi:10.1103/PhysRevLett.48.291}}.
\item
  B. S. Cirel'son, ``Quantum Generalizations of Bell's Inequality,''
  \emph{Letters in Mathematical Physics} \textbf{4}, 93--100 (1980).
  \href{https://doi.org/10.1007/BF00417500}{\nolinkurl{doi:10.1007/BF00417500}}.
\item
  S. Abramsky and A. Brandenburger, ``The Sheaf-Theoretic Structure of
  Non-Locality and Contextuality,'' \emph{New Journal of Physics}
  \textbf{13}, 113036 (2011).
  \href{https://doi.org/10.1088/1367-2630/13/11/113036}{\nolinkurl{doi:10.1088/1367-2630/13/11/113036}}.
\item
  S. Abramsky, S. Mansfield and R. Soares Barbosa, ``The Cohomology of
  Non-Locality and Contextuality,'' \emph{Electronic Proceedings in
  Theoretical Computer Science} \textbf{95}, 1--14 (2012).
  \href{https://doi.org/10.4204/EPTCS.95.1}{\nolinkurl{doi:10.4204/EPTCS.95.1}}.
\item
  N. D. Mermin, ``Simple Unified Form for the Major No-Hidden-Variables
  Theorems,'' \emph{Physical Review Letters} \textbf{65}, 3373--3376
  (1990).
  \href{https://doi.org/10.1103/PhysRevLett.65.3373}{\nolinkurl{doi:10.1103/PhysRevLett.65.3373}}.
\item
  T. Regge, ``General Relativity without Coordinates,'' \emph{Il Nuovo
  Cimento} \textbf{19}, 558--571 (1961).
  \href{https://doi.org/10.1007/BF02733251}{\nolinkurl{doi:10.1007/BF02733251}}.
\item
  Y. Aharonov and D. Bohm, ``Significance of Electromagnetic Potentials
  in the Quantum Theory,'' \emph{Physical Review} \textbf{115}, 485--491
  (1959).
  \href{https://doi.org/10.1103/PhysRev.115.485}{\nolinkurl{doi:10.1103/PhysRev.115.485}}.
\item
  M. V. Berry, ``Quantal Phase Factors Accompanying Adiabatic Changes,''
  \emph{Proceedings of the Royal Society A} \textbf{392}, 45--57 (1984).
  \href{https://doi.org/10.1098/rspa.1984.0023}{\nolinkurl{doi:10.1098/rspa.1984.0023}}.
\item
  A. Baker, ``Linear Forms in the Logarithms of Algebraic Numbers
  (III),'' \emph{Mathematika} \textbf{14}(2), 220--228 (1967).
  \href{https://doi.org/10.1112/S0025579300003843}{\nolinkurl{doi:10.1112/S0025579300003843}}.
\item
  B. Eastin and E. Knill, ``Restrictions on Transversal Encoded Quantum
  Gate Sets,'' \emph{Physical Review Letters} \textbf{102}, 110502
  (2009).
  \href{https://doi.org/10.1103/PhysRevLett.102.110502}{\nolinkurl{doi:10.1103/PhysRevLett.102.110502}}.
\item
  \emph{Native joint operations and a Genesis-to-Bell experiment}, The
  Symbol Game repository, revision
  \texttt{4\allowbreak{}a\allowbreak{}4\allowbreak{}f\allowbreak{}d\allowbreak{}1\allowbreak{}9\allowbreak{}2\allowbreak{}f}
  (2026); included in the code-companion archive as
  \texttt{o\allowbreak{}m\allowbreak{}e\allowbreak{}g\allowbreak{}a\allowbreak{}\_\allowbreak{}e\allowbreak{}n\allowbreak{}g\allowbreak{}i\allowbreak{}n\allowbreak{}e\allowbreak{}/\allowbreak{}c\allowbreak{}r\allowbreak{}a\allowbreak{}t\allowbreak{}e\allowbreak{}s\allowbreak{}\_\allowbreak{}r\allowbreak{}s\allowbreak{}/\allowbreak{}o\allowbreak{}m\allowbreak{}e\allowbreak{}g\allowbreak{}a\allowbreak{}\_\allowbreak{}w\allowbreak{}d\allowbreak{}p\allowbreak{}/\allowbreak{}d\allowbreak{}o\allowbreak{}c\allowbreak{}s\allowbreak{}/\allowbreak{}g\allowbreak{}e\allowbreak{}n\allowbreak{}e\allowbreak{}s\allowbreak{}i\allowbreak{}s\allowbreak{}\_\allowbreak{}b\allowbreak{}e\allowbreak{}l\allowbreak{}l\allowbreak{}\_\allowbreak{}e\allowbreak{}x\allowbreak{}p\allowbreak{}e\allowbreak{}r\allowbreak{}i\allowbreak{}m\allowbreak{}e\allowbreak{}n\allowbreak{}t\allowbreak{}.\allowbreak{}m\allowbreak{}d}.
\item
  A. Baker, ``Linear Forms in the Logarithms of Algebraic Numbers (I),''
  \emph{Mathematika} \textbf{13}, 204--216 (1966).
\end{enumerate}

\end{document}
